% G=a²μ₀ 量纲与数值矛盾的 Lean 4 形式化验证
% 基于塔里芙 (2026) 原文的机器验证报告
% 2026-10-14

\documentclass[12pt,a4paper]{article}
\usepackage{ctex}
\usepackage{amsmath,amssymb,amsthm}
\usepackage{geometry}
\usepackage{hyperref}
\usepackage{booktabs}
\usepackage{listings}
\usepackage{xcolor}
\usepackage{caption}

\geometry{left=2.5cm,right=2.5cm,top=2.5cm,bottom=2.5cm}
\setlength{\emergencystretch}{3em}

\lstset{
  language=Python,
  basicstyle=\small\ttfamily,
  keywordstyle=\color{blue},
  commentstyle=\color{gray},
  numbers=left,
  numberstyle=\tiny\color{gray},
  frame=single,
  breaklines=true
}

\newtheorem{theorem}{定理}
\newtheorem{lemma}[theorem]{引理}
\newtheorem{definition}{定义}

\title{\textbf{$G \approx \alpha^2 \mu_0$ 量纲与数值关系的 Lean 4 辅助验证}}
\author{形式化验证报告 \\ \small 基于塔里芙 (2026) 原文}
\date{2026 年 10 月 14 日}

\begin{document}
\maketitle

\begin{abstract}
塔里芙 (2026) 在《$G=\alpha^2\mu_0$量纲与数值矛盾的完整解决》一文中提出：万有引力常数 $G$ 与精细结构常数平方乘以真空磁导率 $\alpha^2\mu_0$ 之间存在深层联系，其在自然单位制下量纲完全自洽，且数值偏差仅 $0.26\%$。本文以 Lean 4 定理证明器辅助、CODATA 2022 常数、mpmath 高精度重算、\textbf{Arb 认证区间算术}与全族穷举枚举加显式零模型检验，对该文核心断言进行多层次验证。本文的主要结果是\textbf{否定性的}：(1)~原文互为印证的两个公式 $\alpha^2/(c^2\varepsilon_0)$ 与 $\alpha^2\mu_0$ 只是同一恒等式的两种写法，"双重证据"系幻觉（Lean 形式证明）；(2)~SI 中 $G \approx \alpha^2\mu_0$（CODATA 2022 下偏差 $0.2616\%$）是\textbf{单位制依赖}的数值巧合——两数不重合的判定分两层，本文严格区分：\textbf{数学层}，Arb 认证两不确定区间\textbf{严格不相交}，不依赖任何概率假设；\textbf{统计层}，以\textbf{当前} $u(G)$ 计偏差约为 $116\sigma$，该倍数依赖计量能力（若 $u(G)$ 退回 1998 年量级只剩约 $1.7\sigma$），故它是附条件的事实而非定理——在 CGS 高斯制中该关系无对应表达式，在标准自然单位制中 $[G]=[M]^{-2}\neq[\alpha^2\mu_0]=1$，且我们把"$G$ 不进入任何无量纲单项式"形式化为 Lean 定理（不引入质量标度则其指数必为 0）——"量纲矛盾"并未因切换单位制而消解，只能靠重新诠释 $\mu_0$ 量纲的物理假设掩盖；(3)~枢纽构造 $\Lambda = 4\pi^3+\pi^2+\pi = 137.03630 \approx \alpha^{-1} = 137.035999$ 经 $194\,480$ 个低复杂度 $\pi$ 多项式全族穷举确为同族全局最优（次优差 $12$ 倍），但在"靶值与格点无关"的显式零模型下出现同量级最优拟合的概率 $p \approx 0.13$（最保守口径 $0.34$），\textbf{统计不显著}；然而该 $p$ 值随家族边界\textbf{双向漂移}——扩到混合基 $(\pi,e)$ 后不仅 $\Lambda$ 被更好的构造（好 $1.9$ 倍）取代、$p$ 升到 $0.73$，把系数界\textbf{收紧}到 $\lvert a_i\rvert\le 4$ 时同一算法反而给出 $p\approx 0.011$（"显著"）。故 $p$ 不是该构造的内禀属性，本文把它作为"家族依赖的敏感性区间 $0.011$--$0.73$（两端口径不同：$0.0111$ 是同一解析口径收到最小家族的值，$0.729$ 来自 $[100,200]$ 窗口的独立 MC，见 §4.4）报告，而不取单点结论；；(4)~$M^*c^2 = 8.909\times10^{16}$~GeV 与 GUT 能标的"对应"依赖未指明模型的靶值选取，且为 $\Lambda \approx \alpha^{-1}$ 的代数推论。这些数值巧合在性质与命运上都与 Eddington--Dirac 大数纲领同类。全部脚本与机器可读结果已公开发布（代码 MIT、文本与数据 CC BY 4.0），可逐位重跑复现。

\medskip
\noindent\textbf{关键词：} 万有引力常数；精细结构常数；自然单位制；量纲分析；Lean 4；形式化验证
\end{abstract}

\tableofcontents
\newpage

%==============================================================================
\section{引言}
%==============================================================================

万有引力常数 $G$ 是物理学中最古老却最不精确的基本常数之一。其 CODATA 2022 推荐值为
\begin{equation}
  G = 6.67430(15)\times10^{-11}~\text{m}^3\,\text{kg}^{-1}\,\text{s}^{-2},
\end{equation}
相对不确定度约 $2.2\times10^{-5}$，在所有基本常数中精度最低。

另一方面，精细结构常数 $\alpha \approx 1/137.036$ 是量子电动力学的核心无量纲常数，其相对标准不确定度为 $u_r(\alpha) = 1.5\times10^{-10}$。需要特别指出的是，2019 年 SI 重新定义后，真空磁导率 $\mu_0$ 不再是精确定义值 $4\pi\times10^{-7}$~N/A$^2$，而成为由 $\alpha$ 传递不确定度的实验测定量（CODATA 2022：$\mu_0 = 1.256\,637\,061\,27(20)\times10^{-6}$~N/A$^2$；由 $\mu_0 = 2\alpha h/(e^2 c)$ 可知 $u_r(\mu_0) = u_r(\alpha) = 1.5\times10^{-10}$，推导见附录~C）；按 $\varepsilon_0 \equiv 1/(\mu_0 c^2)$ 定义，同理。$c$、$h$、$e$ 在现行 SI 中仍为精确定义值。

塔里芙 (2026) 在公众号文章《$G=\alpha^2\mu_0$量纲与数值矛盾的完整解决》\cite{talihua2026}中注意到一个引人注目的数值巧合：
\begin{equation}\label{eq:main}
  G \approx \alpha^2 \mu_0,
\end{equation}
两者在 SI 单位制中数值相差仅约 $0.26\%$。然而，在标准 SI 量纲分析下，$[G] = \mathsf{L^3 M^{-1} T^{-2}}$ 与 $[\alpha^2\mu_0] = \mathsf{M L T^{-2} I^{-2}}$ 完全不同，构成表观的"量纲矛盾"。

该文提出通过切换到自然单位制（$c = \hbar = 1$）来解决这一矛盾，并引入几何常数 $\Lambda = 4\pi^3+\pi^2+\pi$ 和导出质量尺度 $M^* = M_P/\Lambda$，构建了 $G = \alpha^2/M^{*2}$ 的关系式，声称 $M^*c^2$ 恰好对应大统一能标（GUT, Grand Unified Theory, scale）。

本文的目标是使用 Lean 4 定理证明器\cite{lean4}对该文的全部核心断言进行多层次的机器辅助验证，以判定哪些断言是代数恒等式、哪些是数值巧合、哪些涉及物理假设。

本文的核心贡献有四：(i)~\textbf{揭穿"双重证据"幻觉}——原文互相印证的两个公式实为同一代数恒等式（定理~1，Lean 形式证明）；(ii)~\textbf{确认 $G \approx \alpha^2\mu_0$ 是 SI 单位制特有的数值巧合}——CGS 中该关系无对应表达式、自然单位制中量纲仍不齐次、且不存在承载该"联系"的无量纲不变量（定理~2--3 与 §4.7--4.8，表~\ref{tab:units}）；(iii)~\textbf{对枢纽构造 $\Lambda \approx \alpha^{-1}$ 做事后选择偏差的定量检验}——$194\,480$ 个 $\pi$ 多项式全族穷举显示其为同族最优，但在显式零模型下 $p \approx 0.13$，统计不显著（定理~4），并指出置换检验在此类问题上会退化；同时将其置于 Eddington--Dirac 大数纲领\cite{eddington1936,dirac1937}与 Sakharov 诱导引力\cite{sakharov1967}的思想史坐标中；(iv)~\textbf{提供独立可复现的多层级验证工件}——CODATA 2022 常数下的 mpmath 重算、\textbf{Arb 定向舍入认证区间}\cite{johansson2017}、零模型 Monte Carlo 与枚举脚本全部落盘。

%==============================================================================
\section{物理常数与定义}
%==============================================================================

本文采用 CODATA 2022 推荐值\cite{codata2022}（2019 年 SI 重新定义后的最新一代），定义如下基本常数（表~\ref{tab:constants}）：

\begin{table}[h]
\centering
\caption{基本物理常数（CODATA 2022）}
\label{tab:constants}
\footnotesize
\begin{tabular}{@{}clcll@{}}
\toprule
\textbf{常数} & \textbf{符号} & \textbf{数值} & \textbf{单位} & \textbf{相对不确定度} \\
\midrule
光速 & $c$ & $299\,792\,458$ (精确) & m/s & --- \\
普朗克常数 & $h$ & $6.626\,070\,15\times10^{-34}$ (精确) & J$\cdot$s & --- \\
真空磁导率 & $\mu_0$ & $1.256\,637\,061\,27(20)\times10^{-6}$ & N/A$^2$ & $1.6\times10^{-10}$\textsuperscript{a} \\
真空介电常数 & $\varepsilon_0$ & $8.854\,187\,8188(14)\times10^{-12}$ & F/m & $1.6\times10^{-10}$\textsuperscript{a} \\
精细结构常数 & $\alpha$ & $7.297\,352\,5643(11)\times10^{-3}$ & 无量纲 & $1.5\times10^{-10}$\textsuperscript{a} \\
万有引力常数 & $G$ & $6.674\,30(15)\times10^{-11}$ & m$^3$/(kg$\cdot$s$^2$) & $2.2\times10^{-5}$ \\
元电荷 & $e$ & $1.602\,176\,634\times10^{-19}$ (精确) & C & --- \\
电子质量 & $m_e$ & $9.109\,383\,7139(28)\times10^{-31}$ & kg & $3.1\times10^{-10}$ \\
约化普朗克常数 & $\hbar = h/2\pi$ & $1.054\,571\,817\ldots\times10^{-34}$ & J$\cdot$s & 同 $h$ (精确) \\
普朗克质量 & $M_P$ & $\sqrt{\hbar c/G} = 2.176\,434(24)\times10^{-8}$ & kg & $1.1\times10^{-5}$（$=\tfrac12 u_r(G)$） \\
\bottomrule
\end{tabular}
\end{table}

\noindent\textbf{注\textsuperscript{a}}（$u_r$ 一栏的口径）：本表 $u_r$ 一律\textbf{由本行括号末位反除}：$1.1\times10^{-12}/\alpha=1.51\times10^{-10}$、$2.0\times10^{-16}/\mu_0=1.59\times10^{-10}$、$1.4\times10^{-21}/\varepsilon_0=1.58\times10^{-10}$，于是 $\alpha$ 行是 $1.5\times10^{-10}$、$\mu_0$ 与 $\varepsilon_0$ 两行是 $1.6\times10^{-10}$。但附录~C 的传递链说明：未舍入世界里 $u_r(\alpha)=u_r(\mu_0)=u_r(\varepsilon_0)$ \textbf{严格相等}，所以三个反除值\textbf{不可能同时为真}——它们的差来自 CODATA 把绝对不确定度舍入到两位展示。\par\noindent 更该说的是：这个舍入方向\textbf{对本文不利}。$(20)$ 要成立要求真 $u(\mu_0)\in[1.95,\,2.05)\times10^{-16}$，即公共 $u_r\in[1.552,\,1.631)\times10^{-10}$；$(11)$ 要成立要求真 $u(\alpha)\in[1.05,\,1.15)\times10^{-12}$，即公共 $u_r\in[1.439,\,1.576)\times10^{-10}$。两式交集为 $[1.552,\,1.576)\times10^{-10}$，对应真 $u(\alpha)\ge1.132\times10^{-12}$。本文区间计算取的恰是展示末位 $u(\alpha)=1.1\times10^{-12}$，\textbf{比真值窄 2.9\%--4.5\%}，方向不保守。\par\noindent 为把这一影响钉死，本文另做三档压力测试（同一完全相关 Arb trace，只换 $\alpha$ 端球；工件 \path|src/q5b_alpha_ball_nist_caliber.py| 与 \path|results/q5b_alpha_ball_nist_caliber.json|）：$u(\alpha)=1.10/1.132/1.168\times10^{-12}$ 给 $u_r(\alpha^2\mu_0)=4.52/4.66/4.80\times10^{-10}$（宽度 $+0\%/+2.95\%/+6.15\%$），三档\textbf{全部判不相交}，间距均为 $1.7312553\times10^{-13}=115.42\,u(G)$。口径选择只在 $10^{-21}$ 量级上挪动间距，挪不动“是否重合”的判定。\par\noindent\textbf{注}：2019 SI 重新定义后，$\mu_0$ 与 $\varepsilon_0$ 不再是精确值。$\mu_0$ 的关系式为
\[
  \mu_0 = \frac{2\alpha h}{e^2 c},
\]
其不确定度完全由 $\alpha$ 的实验不确定度传递（$h$、$e$、$c$ 均为精确值）；$\varepsilon_0 \equiv 1/(\mu_0 c^2)$ 为导出量。本文定理~1 的代数恒等式 $c^2\varepsilon_0\mu_0 = 1$ 在定义层面依然成立，不受 SI 重新定义影响；受影响的是把 $\mu_0$ 当作"精确 $4\pi\times10^{-7}$"的表述——本文原始版本（基于 CODATA 2018 框架表述）的这一说法已按 CODATA 2022 修正。

由此导出以下定义：
\begin{definition}[几何常数 $\Lambda$]
\begin{equation}
  \Lambda = 4\pi^3 + \pi^2 + \pi.
\end{equation}
\end{definition}

\begin{definition}[导出质量尺度 $M^*$]
\begin{equation}
  M^* = \frac{M_P}{\Lambda},
\end{equation}
其中 $M_P = \sqrt{\hbar c / G} = 2.176\,434(24)\times10^{-8}$~kg（对应 $M_Pc^2 = 1.2209\times10^{19}$~GeV）为\textbf{普朗克质量}。注意本文不使用"约化普朗克质量"约定：约化版为 $\overline{M}_P = \sqrt{\hbar c/(8\pi G)} = M_P/\sqrt{8\pi} = 2.4353\times10^{18}$~GeV$/c^2$，两者相差 $\sqrt{8\pi}\approx5.01$ 倍，混用会把 $M^*c^2$ 挪动同样的倍数。原文与本节所有 $M_P$ 均为非约化定义，自然单位制下 $G_{\rm nat} = 1/M_P^2 = 6.709\times10^{-39}$~GeV$^{-2}$（与文献值 $6.7088\times10^{-39}$~GeV$^{-2}$ 一致，可用作约定自检）。
\end{definition}

%==============================================================================
\section{方法：Lean 4 形式化验证}
%==============================================================================

本文使用 Lean 4.33.1 定理证明器\cite{lean4}辅助，结合 Python mpmath 高精度计算库\cite{mpmath}，对全部断言进行多层次验证。需要明确指出的是，各断言的验证深度不同：

\begin{enumerate}
  \item \textbf{形式证明}（Lean 4）：两类断言已机械化。(a) 代数恒等式（定理~1）由 \texttt{by rfl} 从定义直接推出；(b) \textbf{量纲格的不变量分类}（§4.7(四)：无量纲单项式的完整刻画、"无质量标度则 $G$ 不进入任何无量纲乘积"、$[G]\neq[\mu_0]$）与 \textbf{Arb 证书的回放}（§4.8）在 \texttt{G\_dimension\_lattice.lean} 中证明，全部经 \texttt{\#print axioms} 核验仅依赖 Lean 的 \texttt{propext}/\texttt{Classical.choice}/\texttt{Quot.sound} 三公理，无 \texttt{sorryAx}。
  \item \textbf{数值验证}（Lean 4 \texttt{Float} + Python mpmath）：对涉及浮点计算的断言（定理~2、4、5、6），使用 IEEE 754 双精度浮点数（Lean \texttt{Float} 类型）进行计算，并用 Python mpmath（dps=30/50）独立重算交叉验证。需注意，数值验证不构成数学证明——它只验证了特定数值输入下的特定精度输出。
  \item \textbf{解析推导}（未形式化）：自然单位制下 $\mu_0$ 量纲重诠释的论证（定理~3 步骤 4--5）仍以解析推导给出。其\textbf{前置事实}（单项式层面的不变量分类）现已形式化，但"赋 $\mu_0$ 以 $[M]^{-2}$ 是否为物理上可接受的假设"属物理论证，不在形式化范围内。
  \item \textbf{认证区间}（Arb 球算术\cite{johansson2017}，§4.8）：以 CODATA 2022 不确定度为输入球，逐步定向舍入，证明 $\alpha^2\mu_0$ 与 $G$ 的区间不相交、$\Lambda \neq \alpha^{-1}$、$\alpha^2\Lambda^2 \neq 1$。mpmath \texttt{mpi}（误差传播包络，非逐舍入认证）仅作交叉核对保留。
  \item \textbf{穷举枚举 + 零模型检验}（Python，§4.4）：对 $194\,480$ 个低复杂度 $\pi$ 多项式全族穷举，在"靶值与格点无关"的显式零模型下给出 $p$ 值与顺序统计量，并把搜索类扩到 $\pi$ 的有理函数、混合基 $(\pi,e)$ 与半整数幂以检验家族稳健性。
\end{enumerate}

阈值选择说明：各数值检查的判定阈值（如 $0.27\%$、$0.01$、$0.002$）只是\textbf{通过性判据}，不构成证据强度——其中 $0.27\%$ 与实际偏差 $0.2616\%$ 仅差 $3\%$，正是紧贴数据点。本文的实质结论不依赖这些阈值，而由 §4.4 的显式零模型与 §4.8 的认证区间给出。

全部源代码与机器可读结果已公开发布于 \path|https://github.com/jackyfan01/g-alpha-mu0-artifacts|（代码 MIT、论文文本与结果 JSON 及证明工件 CC BY 4.0；仓库内含 \path|paper/| 目录，收录本文 \path|.tex|/\path|.md|/\path|.pdf| 三形态；Zenodo DOI 于首个 release 后写入本节）；发布前的本地同名目录为 \path|G_alpha_mu0_repo/|（含 \path|README.md|、\path|environment.yml|、\path|requirements.txt|、\path|lakefile.lean|、\path|lean-toolchain|）：Lean 侧为 \path|src\G_alpha_squared_mu_zero.lean|（定理~1--6 的数值检查与量纲注释）与 \path|src\G_dimension_lattice.lean|（量纲格与证书回放的形式化）；Python 侧为 \path|src\verify_numerical.py|（全量重算）、\path|src\R3_lambda_search.py|（全族穷举）、\path|src\q2_lambda_null_model.py|（零模型与 $p$ 值）、\path|src\q3b_family_extension.py|（家族稳健性）、\path|src\q3_arb_certified.py|（Arb 认证区间）、\path|src\R4_interval_arithmetic.py|（mpmath 交叉核对）、\path|src\q5_correlated_and_spacing.py|（完全相关证书与间距感知精确 $p$）与 \path|src\q5b_alpha_ball_nist_caliber.py|（$\alpha$ 端球口径的三档压力测试），机器可读结果在 \path|results\*.json|。运行环境：Python 3.13.7、mpmath 1.3.0、numpy 2.4.3、python-flint 0.9.0（Arb 2.x）、Lean 4.33.1；随机实验固定种子（\texttt{q2}: 20261014，\texttt{q3b}: 20261015），重跑逐位复现。

%==============================================================================
\section{结果}
%==============================================================================

\subsection{定理 1：代数恒等式 $\alpha^2/(c^2\varepsilon_0) = \alpha^2\mu_0$}

\begin{theorem}[代数等价性]
由真空介电常数的定义 $\varepsilon_0 \stackrel{\text{def}}{=} 1/(\mu_0 c^2)$，有
\begin{equation}
  \frac{\alpha^2}{c^2 \varepsilon_0} = \alpha^2 \mu_0.
\end{equation}
\end{theorem}

\begin{proof}
由定义，$c^2 \varepsilon_0 = c^2 \cdot \frac{1}{\mu_0 c^2} = \frac{1}{\mu_0}$。因此
\[
  \frac{\alpha^2}{c^2 \varepsilon_0} = \alpha^2 \cdot \mu_0. \qed
\]
\end{proof}

\noindent\textbf{Lean 形式证明}：\texttt{成立}（\texttt{G\_dimension\_lattice.lean} 中的 \texttt{thm1\_monoid}，\texttt{by rfl}，\texttt{\#print axioms} 显示\textbf{不依赖任何公理}）。诚实限定：该证明在\textbf{单项式的自由阿贝尔群}（以 $(\alpha,\mu_0,c)$ 为基的整数指数向量）中给出——这正是断言所在的层面；在任意域上的证明需 Mathlib 的 \texttt{field\_simp}，本文不依赖 Mathlib 故未做。此式为由定义保证的代数恒等式，与物理无关。

\medskip

\textbf{物理意义}：原文将 $G = \alpha^2/(c^2\varepsilon_0)$ 与 $G = \alpha^2\mu_0$ 写成两个独立公式，读起来像互相印证。但由定理~1，两者实为同一关系的两种写法，不构成"双重验证"。

\subsection{定理 2：SI 中 $G \approx \alpha^2\mu_0$（数值验证）}

\begin{theorem}[数值近似]
采用 CODATA 2022 常数，在 SI 单位制中，
\begin{equation}
  \alpha^2 \mu_0 = 6.691\,762\,6\times10^{-11}~\text{N\,A}^{-2}
    \quad(\text{量纲 }\mathsf{M\,L\,T^{-2}I^{-2}}),
\end{equation}
而 $G = 6.674\,30(15)\times10^{-11}~\text{m}^3\,\text{kg}^{-1}\,\text{s}^{-2}$
（量纲 $\mathsf{L^{3}M^{-1}T^{-2}}$）。两者\textbf{数值}的相对偏差为 $0.2616\%$；
但它们的量纲不同，故该比较只在 SI 数值约定下有意义（表~\ref{tab:units}、§4.7）。
\end{theorem}

\noindent\textbf{数值检查}：\texttt{PASS}（偏差 $0.2616\% < 0.27\%$；Lean \texttt{Float} 与 mpmath 双通道）。此为浮点计算结果，\textbf{不是}形式证明；其严格版本见 §4.8 的 Arb 认证区间。

\medskip

\textbf{不确定度传播}：$\alpha^2\mu_0$ 的相对不确定度\textbf{不能}按相互独立的 $A$、$B$ 类分量做平方和开根。因为在现行 SI 中 $\mu_0 = 2\alpha h/(e^2 c)$ 而 $h$、$e$、$c$ 均为精确定义值，$\mu_0$ 是不含任何其他输入的 $\alpha$ 的确定函数，$u_r(\mu_0) = u_r(\alpha)$ 且与 $\alpha$ \textbf{完全正相关}（附录~C 给出推导）。于是线性传播（全相关）给出
\[
  u_r(\alpha^2\mu_0) = 2u_r(\alpha) + u_r(\mu_0) = 3u_r(\alpha) \approx 4.5\times10^{-10},
\]
即 $\alpha^2\mu_0$ 的有效数字精度约 9--10 位；若误按独立分量合成会得到 $\sqrt{5}\,u_r(\alpha)\approx 3.4\times10^{-10}$，那是\textbf{低估}。而 $G$ 的相对不确定度为 $2.2\times10^{-5}$，比前者大 $4.97\times10^{4}$ 倍（接近五个数量级）。因此 $0.2616\%$ 的偏差完全由 $G$ 一侧主导，$\mu_0$ 从"精确值"改为"实验测定值"（CODATA 2022）对该结论\textbf{无任何影响}——偏差量级 $10^{-3}$ 远大于双方不确定度之和（$\sim 2.2\times10^{-5}$，即 §4.8 的 $116\sigma$）。§4.8 的 Arb 认证区间进一步给出两区间不相交的机器可复查证明。

\medskip

\textbf{分析}：此数值关系是单位制依赖的。在 SI 中，$G$ 和 $\alpha^2\mu_0$ 的量纲不同（分别为 $\mathsf{L^3M^{-1}T^{-2}}$ 和 $\mathsf{MLT^{-2}I^{-2}}$），因此"数值相近"会随单位制改变而消失。例如在 CGS 高斯制中，两者的数值不再接近（详见 §4.7 的无量纲不变量分析）。这正是原文承认"量纲匹配问题没解决"的根源。

\subsection{定理 3：自然单位制下量纲自洽}

\begin{theorem}[量纲分析]
在自然单位制（$c = \hbar = 1$，采用 Heaviside--Lorentz（HL）约定\cite{peskin1995}）下，
\begin{equation}
  [G] = [M]^{-2}.
\end{equation}
若要求 $G \approx \alpha^2\mu_0$ 在自然单位制中成为量纲齐次的关系式，则 $\mu_0$ 须被赋予量纲 $[\mu_0] = [M]^{-2}$，而非标准自然单位制中的无量纲。
\end{theorem}

\begin{proof}
在自然单位制中，$[c] = [\hbar] = 1$。由 $[\hbar] = [M][L]^2[T]^{-1} = 1$ 和 $[c] = [L][T]^{-1} = 1$ 得 $[L] = [T] = [M]^{-1}$。

\textbf{步骤 1}：$[G]$ 的独立推导。SI 中 $[G] = \mathsf{L^3 M^{-1} T^{-2}}$，代入 $[L]=[T]=[M]^{-1}$ 得
\[
  [G] = [M]^{-3} \cdot [M]^{-1} \cdot [M]^{2} = [M]^{-2}.
\]

\textbf{步骤 2}：$[\alpha^2]$ 的独立推导。精细结构常数 $\alpha$ 始终无量纲，故 $[\alpha^2] = 1$。

\textbf{步骤 3}：$[\mu_0]$ 在标准自然单位制中的推导。在 Heaviside--Lorentz 自然单位制中，$\alpha = e^2/(4\pi)$，故 $[e] = 1$。由 $\alpha = e^2/(4\pi\varepsilon_0)$ 且 $[\alpha]=[e]=1$，得 $[\varepsilon_0] = 1$。再由 $c^2 = 1/(\mu_0\varepsilon_0)$ 且 $[c]=1$，得 $[\mu_0] = 1$（无量纲）。

\textbf{步骤 4}：量纲矛盾的重新表述。由步骤 2--3，$[\alpha^2\mu_0] = 1 \cdot 1 = 1$。但由步骤 1，$[G] = [M]^{-2} \neq 1$。因此，在标准自然单位制下，$G = \alpha^2\mu_0$ \textbf{并不量纲齐次}——量纲矛盾并未因切换到自然单位制而消失，而是从 SI 中的 $[G] \neq [\alpha^2\mu_0]$（两种不同的复合量纲）转化为自然单位制中的 $[M]^{-2} \neq 1$（质量量纲 vs 无量纲）。

\textbf{步骤 5}：重新诠释 $\mu_0$ 的量纲。若要将 $G \approx \alpha^2\mu_0$ 理解为自然单位制中的量纲齐次关系，则须重新赋予 $\mu_0$ 以量纲 $[\mu_0] = [M]^{-2}$。这等价于在自然单位制中将 $\mu_0$ 视为一个具有质量量纲的导出量（而非标准电磁学中的无量纲常数），其物理含义是 $\mu_0$ 的数值与引力耦合常数 $G$ 通过 $\mu_0 \sim G/\alpha^2$ 关联。这一赋量纲操作在标准场论中\textbf{没有依据}（见 §5.2 的讨论）：它是一步物理假设，而非单位制切换的自动后果。
\end{proof}

\noindent\textbf{解析推导}（未形式化）：量纲分析在 Lean 代码注释中给出，Lean 4 未对此作形式验证。步骤 1--3 为标准量纲分析，步骤 4--5 为本文的新分析，揭示了原文"量纲矛盾已解决"的断言需要更精确的表述。三种单位约定下的量纲对照见表~\ref{tab:units}。

\subsection{定理 4：$\Lambda \approx 137 \approx \alpha^{-1}$}

\begin{theorem}[$\Lambda$ 的数值]
\begin{align}
  \Lambda &= 4\pi^3 + \pi^2 + \pi = 137.03630, \\
  \alpha^{-1} &= 137.035999178 \quad (\text{CODATA 2022}), \\
  |\Lambda - 137| &= 0.03630, \\
  |\Lambda - \alpha^{-1}| &= 3.046\times10^{-4}.
\end{align}
\end{theorem}

\noindent\textbf{数值检查}：\texttt{PASS}（$|\Lambda - 137| < 0.1$ 且 $|\Lambda - \alpha^{-1}| < 0.01$；Arb 认证区间见 §4.8）。

\medskip

\textbf{分析}：$\Lambda = 4\pi^3+\pi^2+\pi$ 是一个纯几何常数（仅依赖 $\pi$），其数值 $137.03630$ 与 $\alpha^{-1} = 137.035999$ 的吻合度达 $3.0\times10^{-4}$。把 $\Lambda$ 代入 $M^*=M_P/\Lambda$ 可得：$\Lambda\approx\alpha^{-1}$ 与 $\alpha^2/M^{*2}\approx G_{\rm nat}$ 是同一近似的两种写法（推导见 §4.6 定理~6），因此这里不存在第二个独立巧合。

\medskip

\textbf{事后选择偏差检验（系统枚举）}：$\Lambda$ 是"先知道 $\alpha^{-1}\approx 137$、再挑选 $\pi$ 多项式逼近它"的事后构造，还是低复杂度 $\pi$ 多项式家族中罕见的优拟合？本文对系数 $a_i \in [-10, 10]$、阶数 $\le 3$ 的\textbf{全体}整数 $\pi$ 多项式
\[
  P(\pi) = a_0 + a_1\pi + a_2\pi^2 + a_3\pi^3
\]
进行了穷举（$21^4 - 1 = 194\,480$ 个候选，排除全零），逐一计算 $|P(\pi) - \alpha^{-1}|$（CODATA 2022，mpmath dps=50，脚本 \texttt{src/R3\_lambda\_search.py}，结果 \texttt{results/lambda\_enumeration.json}）。结论分两层：

\begin{enumerate}
  \item \textbf{在全族 $194\,480$ 个候选中，$\Lambda=(0,1,1,4)$ 是唯一的全局最优}（排名 0），偏差 $3.046\times10^{-4}$；次优候选 $(-4,9,2,3)$ 偏差 $3.626\times10^{-3}$，比 $\Lambda$ 差 $12$ 倍；第三 $(4,-7,0,5)$ 偏差 $4.236\times10^{-3}$。偏差 $<10^{-3}$ 的仅 $\Lambda$ 一个，$<10^{-2}$ 的仅 3 个，$<0.1$ 的 68 个。就"同族比较"而言，$\Lambda$ 确实是构造者能找到的最佳拟合——这解释了它为何被选中。
  \item \textbf{但"最佳"不等于"统计显著"——显式零模型（五种口径）}。构造零假设 $H_0$：靶值 $T$ 与该格点的取值集合无关（$T$ 在集合覆盖的区间 $[-450.17, 450.17]$ 上均匀分布），统计量为 $D_{\min}(T)=\min_i |V_i - T|$，其中 $V_i$ 遍历家族取值、$R=V_{\max}-V_{\min}=900.35$。以 $\alpha^{-1}$ 的观测偏差 $d_{\text{obs}} = 3.0460\times10^{-4}$ 为临界值（脚本 \texttt{src/q2\_lambda\_null\_model.py}，种子 20261014，结果 \texttt{results/q2\_lambda\_null\_results.json}），四种口径给出：
  \begin{center}
  \footnotesize
  \begin{tabular}{@{}lll@{}}
  \toprule
  \textbf{口径} & \textbf{$p$ 值} & \textbf{说明} \\
  \midrule
  随机靶 Monte Carlo（$4\times10^5$ 次，精确最近邻） & $0.127$ & $d_{\text{obs}}$ 落在 $D_{\min}$ 分布的第 $12.7$ 百分位 \\
  解析二项式 $1-(1-2d/R)^N$ & $0.123$ & 与 MC 一致 \\
  Bonferroni 上界 $N\cdot 2d/R$ & $0.132$ & 多重检验校正后的保守上界 \\
  限制在物理合理窗口 $T\in[100,200]$ & $0.173$ & 回应"家族定义过宽"的变体 \\
  条件于"$0.1$ 邻域内已有 68 个候选" & $0.340$ & 最保守口径（本文初版采用） \\
  \bottomrule
  \end{tabular}
  \end{center}
  全部口径均不低于 $0.05$（$p=0.127$ 的 MC 标准误约 $5.3\times10^{-4}$），即\textbf{无法拒绝}$H_0$。必须把这一结论的方向说准：它\textbf{不是}"已证明 $\Lambda \approx \alpha^{-1}$ 纯属巧合"，而是"该吻合未提供足以区分它与随机挑选的证据"——$1-p$ 不是"$H_0$ 为假"的概率，零模型本身（$T$ 在家族值域上均匀、家族边界取 $\lvert a_i\rvert\le10$、$\le3$ 阶）也是人为指定的，换成 $T\sim U[100,200]$ 口径（家族值域取 $[100,200]$）$p$ 变为 $0.175$，即下文家族稳健性表第一行。\par\medskip\noindent\textbf{同一零模型的解析精确值（去掉 Monte Carlo 噪声）}。这里先纠正一个容易说错的前提：均匀随机靶\textbf{并没有}忽略格点覆盖的不均匀性——靶落进第 $i$ 个 Voronoi 胞的概率恰为 $w_i/R$（$w_i$ 是胞宽），均匀口径\textbf{本身就}是按胞宽加权；而胞内任一点到本胞站点的距离，就是它到\textbf{最近}格点的距离。所以“按胞宽加权抽样”与“均匀随机靶”是同一个模型，其被估量可直接写成最近邻覆盖测度\[ p_{\mathrm{exact}} \;=\; \frac{\bigl|\,\bigcup_i [V_i-d_{\mathrm{obs}},\,V_i+d_{\mathrm{obs}}] \cap [V_{\min},V_{\max}]\,\bigr|}{V_{\max}-V_{\min}}, \]区间合并后逐段求长，公式层面无任何近似（工件用 float64 求值，另以 mpmath dps=60 复核五档，与工件在小数点后六位全同，最大偏差 $3.9\times10^{-14}$）。实测（$d_{\mathrm{obs}} = 3.04598\times10^{-4}$）：\begin{center}\footnotesize\begin{tabular}{@{}lrrrr@{}}\toprule\textbf{家族（$\pi$ 的 $\le3$ 次整系数多项式）} & $N$ & 值域 $R$ & $p_{\mathrm{exact}}$ & $p_{\mathrm{Bonf}}$ \\\midrule$\lvert a_i\rvert\le 4$（= $\Lambda$ 自身的最大系数） & 6\,560 & $360.14$ & $0.0111$ & $0.0111$ \\$\lvert a_i\rvert\le 5$ & 14\,640 & $450.17$ & $0.0198$ & $0.0198$ \\$\lvert a_i\rvert\le 6$ & 28\,560 & $540.21$ & $0.0322$ & $0.0322$ \\$\lvert a_i\rvert\le 8$ & 83\,520 & $720.28$ & $0.0698$ & $0.0706$ \\$\lvert a_i\rvert\le 10$（正文主口径） & 194\,480 & $900.35$ & $\mathbf{0.1270}$ & $0.1316$ \\\bottomrule\end{tabular}\end{center}\noindent 主口径下 $p_{\mathrm{exact}} = 0.126988$ 与独立 Monte Carlo 的 $0.12708$（$4\times10^5$ 次）在\textbf{小数点后三位}一致（同为 $0.127$），差 $9.2\times10^{-5}$，仅为主口径 MC 标准误 $5.3\times10^{-4}$ 的 $0.17$ 倍——\textbf{它们本是同一被估量的两种算法}，相符是应当的，不相符才说明其中必有一处算错。另有三点必须同时说清：(i) 本表的 $p$ 随家族规模从 $0.0111$ 单调升到 $0.1270$，把上文“缩族/扩族”的双向漂移变成一条可读的单调曲线；但正文另处引用的混合基 $(\pi,e)$ 值 $0.729$ \textbf{不属于这条曲线}——它是另一口径（靶均匀于 $[100,200]$ 的 MC，$2\times10^5$ 次，种子 20261015），不能与 $p_{\mathrm{exact}}$ 串联比较。同一 $[100,200]$ 口径本文有两次独立 MC：$0.17311$（$4\times10^5$，种子 20261014）与 $0.17511$（$2\times10^5$，种子 20261015），差 $2.0\times10^{-3}\approx1.9$ 倍合并标准误，属抽样差，两处引用时须注明是哪一次。(ii) 逐胞求长时\textbf{不能}用对称的 $\min(2d,w_i)/R$：本家族有 $17.15\%$ 的相邻间距小于 $2d_{\mathrm{obs}}$，该近似给 $0.1310$，比真值大 $3.2\%$；正确式是 $\sum_i\bigl(\min(d,L_i)+\min(d,R_i)\bigr)/R$（$L_i,R_i$ 为站点到胞左右边界的距离），它与覆盖测度在本例相符到 $1.3\times10^{-15}$。(iii) 把“家族本身也是挑的”计入多重比较：对 $K=5$ 个口径取最小 $p$ 后 Bonferroni 得 $0.0110949\times5=0.0555$。但这五个口径是\textbf{嵌套}的（$\Lambda$ 属于每一个，$B\ge4$），最小 $p$ 必然落在最小家族，$K=5$ 本身又是人为选的，故 $0.0555$ \textbf{只能当敏感性指标，不能当频率学 $p$ 值}（工件 \texttt{q5\_correlated\_and\_spacing.json} 的 \texttt{cross\_family.note} 亦作此声明）。它的正当用途只有一个：说明\textbf{单看任何一个口径的 $p$ 都是在做选择}。因此可辩护的读法只有一条：在给定家族与给定零模型下，$\Lambda$ 的"最优性"不足以充当 $\alpha^{-1}$ 由纯几何决定的证据；它是事后选择偏差（a posteriori selection bias）的典型形态。顺序统计量还显示，随机靶下 $D_{\min}$ 的中位数是 $1.53\times10^{-3}$、5\% 分位是 $1.16\times10^{-4}$——$\Lambda$ 的拟合优度只是"比典型好一点"，远非异常。
  \item \textbf{一个方法论副产物：系数空间置换检验在此类问题上退化}。对系数空间做 permutation/bootstrap 是这类拟合问题的常规稳健性检验，本文也照此尝试。但 $\Lambda$ 是全格点的\textbf{全局最优}，故任何由格点成员构成的子集，其最佳拟合距离恒 $\ge d_{\text{obs}}$，等号成立当且仅当子集恰好包含 $\Lambda$。实测（200 次重抽）确认：所有 $D_{\min} \ge d_{\text{obs}}$，$P(D_{\min} \le d_{\text{obs}}) = 0.585 \approx 1-e^{-1}$——这正是"$N$ 次有放回抽样抽中特定成员"的概率，与拟合优劣\textbf{完全无关}。因此本问题的合法零模型只能是上面的随机靶口径；对"事后挑选"型问题，置换检验若以被选对象自身为参照，会退化为选择概率的度量。
  \item \textbf{家族外稳健性：扩大搜索类后 $\Lambda$ 连"最优"都不再是}。上面的结论依赖家族定义（整数系数、$\pi$ 的 $\le 3$ 次多项式）。为检验这一家族依赖性，把搜索类换为四个不同家族，统一以 $T \sim U[100,200]$ 为随机靶（脚本 \texttt{src/q3b\_family\_extension.py}，种子 20261015）：

  \begin{center}
  \footnotesize
  \scriptsize
  \begin{tabular}{@{}lrrrl@{}}
  \toprule
  \textbf{家族} & $N$ & \textbf{最佳距离} & \textbf{优于 $\Lambda$?} & $p$ \\
  \midrule
  $\pi$ 多项式，$\lvert a_i\rvert\le10$，$\le 3$ 阶 & 194\,480 & $3.046\times10^{-4}$ & —（即 $\Lambda$） & $0.175$ \\
  $\pi$ 的有理函数 $\le$ 2 阶，$\lvert a_i\rvert\le4$ & 59\,049 & $3.6\times10^{1}$ & 否（家族在该处过稀） & $0.371$ \\
  混合基 $(\pi,e)$，$\lvert a_i\rvert\le5$，6 系数 & 1\,771\,560 & $\mathbf{1.600\times10^{-4}}$ & \textbf{是（好 1.9 倍）} & $\mathbf{0.729}$ \\
  $\pi$ 的半整数幂 $\{1,\pi^{3/2},\pi^{5/2}\}$ & 9\,260 & $1.374\times10^{-2}$ & 否 & $0.355$ \\
  \bottomrule
  \end{tabular}
  \end{center}

  关键结果在第三行：把基从 $\pi$ 扩到 $\{\pi, e\}$、系数从 4 个扩到 6 个后，出现
  \[
    -1 - 2\pi + 5\pi^2 + 3\pi^3 - 2e + e^2 = 137.03615918,
    \qquad d = 1.600\times10^{-4},
  \]
  比 $\Lambda$ 好 $1.9$ 倍——$\Lambda$ \textbf{不再是全局最优}。而该家族下的偶然概率高达 $p = 0.729$：家族一大，"找到更好的拟合"几乎是必然事件。\textbf{但本文必须同时报告反方向的敏感性检验——缩小家族会让"$p\approx 0.13$ 不显著"翻转。}把系数界从 $\lvert a_i\rvert\le 10$ 分别收紧到 $\le 5$（$N=14\,640$）与 $\le 4$（$N=6\,560$，恰为 $\Lambda=(0,1,1,4)$ 自身的最大系数），在同一 Bonferroni 口径下$p$ 由 $0.132$ 降为 $0.0198$ 与 $0.0111$；解析式 $1-(1-2d/R)^N$ 给出 $0.0196$ 与 $0.0110$（本人 2026-10-16 重跑 $N$ 与值域 $R$ 实算，非引用他人结果）。也就是说：按这套算法，$\Lambda$ 在\textbf{更低复杂度界下反而是"显著"的}，而 $p$ 值随家族边界\textbf{双向漂移}（扩族升到 $0.17$--$0.73$，缩族降到 $0.011$）。因此"$p\approx 0.13$"不是该构造的内禀属性，任何单边引用都会误导。这从反面印证了本文的定性：$\Lambda$ 的"同族最优"是\textbf{家族规模}的产物，不是 $\alpha^{-1}$ 的几何必然。任何以"$\pi$ 的某个组合凑出 137"为卖点的构造，都必须在扩大家族后重新检验——而扩大之后总有更"好"的构造出现，这正是数值命理学（numerology）无法自证的结构性原因。
\end{enumerate}

\subsection{定理 5：$M^*c^2$ 与大统一能标的对应（模型依赖）}

\begin{theorem}[GUT 能标对应]
\begin{equation}
  M^* c^2 = \frac{M_P c^2}{\Lambda} = 8.909\times10^{16}~\text{GeV},
\end{equation}
该数值与\textbf{特定取法}的"GUT 能标" $E_\text{GUT} \approx 8.9\times10^{16}$~GeV 偏差 $< 1\%$。但"GUT 能标"并非单一数值，而是各具体统一模型重整化群拟合的产物；GUT 即 Grand Unified Theory（大统一理论），其能标不是可观测量，换模型就会移动（见表~\ref{tab:gut}）。
\end{theorem}

\begin{table}[h]
\centering
\caption{$M^*c^2 = 8.909\times10^{16}$~GeV 与若干典型 GUT 模型能标的比较（模型能标为量级估计）}
\label{tab:gut}
\footnotesize
\begin{tabular}{@{}lcc@{}}
\toprule
\textbf{统一模型} & \textbf{GUT 能标 (GeV)} & \textbf{与 $M^*c^2$ 的比} \\
\midrule
最小非超对称 SU(5) & $\sim 10^{15}$ & $M^*c^2$ 高约 $89$ 倍 \\
最小 SO(10) & $\sim 10^{16}$ & $M^*c^2$ 高约 $8.9$ 倍 \\
SUSY SU(5) / MSSM 规范耦合统一 & $\sim 2\times10^{16}$ & $M^*c^2$ 高约 $4.5$ 倍 \\
质子衰变实验下限 & $M_X \gtrsim 3\times10^{15}$（$X$ 玻色子质量） & — \\
原文所引靶值 & $8.9\times10^{16}$ & 偏差 $< 1\%$ \\
\bottomrule
\end{tabular}
\end{table}

\noindent 表中模型能标为文献中的常引量级估计（综述见 \cite{nath2007}）；MSSM = Minimal Supersymmetric Standard Model（最小超对称标准模型），$X$ 玻色子为重子数破坏的规范玻色子。下限来自质子衰变实验搜寻，故为间接界限而非直接测量。\par\medskip

\noindent\textbf{数值检查}：\texttt{PASS}（数值本身 $8.909\times10^{16}$~GeV 无误，Arb 认证区间 $[8.9091,8.9093]\times10^{16}$~GeV；但"对应"的靶值选取是模型依赖的，见 §5.5）。

\subsection{定理 6：$G \approx \alpha^2/M^{*2}$ 的自洽性}

\begin{theorem}[核心公式自洽性]
在自然单位制中，
\begin{equation}
  \frac{G}{\alpha^2/M^{*2}} = \frac{1}{\alpha^2\Lambda^2} = 0.999996,
\end{equation}
其中 $\alpha^2\Lambda^2 = 1.000004$。即 $G \approx \alpha^2/M^{*2}$，相对偏差仅 $4.45\times10^{-6}$（$=1-1/\alpha^2\Lambda^2$）。
\end{theorem}

\begin{proof}
在自然单位制中，$G = 1/M_P^2$（$M_P$ 为普朗克质量，非约化版），且 $M^* = M_P/\Lambda$，故
\[
  \frac{\alpha^2}{M^{*2}} = \frac{\alpha^2}{M_P^2/\Lambda^2} = \alpha^2\Lambda^2 \cdot \frac{1}{M_P^2} = \alpha^2\Lambda^2 \cdot G.
\]
因此 $G / (\alpha^2/M^{*2}) = 1/(\alpha^2\Lambda^2)$。

由定理~4，$\Lambda \approx \alpha^{-1}$，故 $\alpha^2\Lambda^2 \approx \alpha^2 \cdot \alpha^{-2} = 1$。

精确计算（CODATA 2022）：$\alpha^2\Lambda^2 = (7.2973525643\times10^{-3})^2 \times (137.03630)^2 = 1.0000044$。
\end{proof}

\noindent\textbf{数值检查}：\texttt{PASS}（$|\alpha^2\Lambda^2 - 1| = 4.4\times10^{-6} < 0.002$）。Arb 认证区间给出 $\alpha^2\Lambda^2 = 1.000004446 \pm 7.8\times10^{-10}$，\textbf{不包含} 1（§4.8），故该偏差不是数值误差造成的表观现象。

\subsection{无量纲不变量分析与单位制无关性}

本节回答一个必须正面处理的问题：$G \approx \alpha^2\mu_0$ 是否在更换单位制后仍然成立，是否存在单位制无关的对应无量纲关系。本节给出定量回答。为避免全文在三种单位约定间来回跳转，先给出一张对照表。

\begin{table}[h]
\centering
\caption{三种单位约定下相关量的量纲与数值}
\label{tab:units}
\scriptsize
\begin{tabular}{@{}llll@{}}
\toprule
\textbf{量} & \textbf{SI} & \textbf{Gaussian CGS} & \textbf{HL 自然单位制} \\
\midrule
$c$ & $\mathsf{L\,T^{-1}}$，$2.998\times10^{8}$ m/s（精确） & $\mathsf{L\,T^{-1}}$，$2.998\times10^{10}$ cm/s & $1$ \\
$\hbar$ & $\mathsf{M\,L^2\,T^{-1}}$，$1.055\times10^{-34}$ J$\cdot$s（精确） & 同量纲，$1.055\times10^{-27}$ erg$\cdot$s & $1$ \\
$e$ & $\mathsf{I\,T}$，$1.602\times10^{-19}$ C（精确） & $\mathsf{M^{1/2}L^{3/2}T^{-1}}$，$4.803\times10^{-10}$ statC & 无量纲，$\alpha=e^2/4\pi$ \\
$\alpha$ & $1$，$7.2973525643(11)\times10^{-3}$ & $1$，同值 & $1$，同值 \\
$\varepsilon_0$ & $\mathsf{M^{-1}L^{-3}T^{4}I^{2}}$，$8.854\times10^{-12}$ F/m & $1$（被吸收） & $1$ \\
$\mu_0$ & $\mathsf{M\,L\,T^{-2}I^{-2}}$，$1.2566\times10^{-6}$ N/A$^2$ & $1$（被吸收） & $1$ \\
$G$ & $\mathsf{L^{3}M^{-1}T^{-2}}$，$6.6743\times10^{-11}$ & $\mathsf{L^{3}M^{-1}T^{-2}}$，$6.6743\times10^{-8}$ & $\mathsf{[M]^{-2}}$，$6.709\times10^{-39}$ GeV$^{-2}$ \\
$\alpha^2\mu_0$ & $\mathsf{M\,L\,T^{-2}I^{-2}}$，$6.6918\times10^{-11}$ & $1$（无量纲），$5.325\times10^{-5}$ & $1$（无量纲），$5.325\times10^{-5}$ \\
\bottomrule
\end{tabular}
\end{table}

\noindent 表~\ref{tab:units} 最后一行是全部问题的核心：$\alpha^2\mu_0$ 的\textbf{数值}在 SI 里恰好与 $G$ 同量级（$6.69$ vs $6.67$，且都带 $10^{-11}$ 量级），但这纯属巧合——它的\textbf{量纲} $\mathsf{M\,L\,T^{-2}I^{-2}}$ 与 $G$ 的 $\mathsf{L^{3}M^{-1}T^{-2}}$ 不同，换到 CGS 或自然单位制后两者立刻分处不同量纲（$5.3\times10^{-5}$ 无量纲 vs $10^{-8}$ 有量纲 / $[M]^{-2}$），比较失去意义。

\textbf{（一）CGS 高斯制中的检验}。CGS 高斯制中不存在独立的 $\mu_0$（其值被吸收进电荷单位定义，$\mu_0^{\text{Gauss}} = 1$ 无量纲），$G = 6.6743\times10^{-8}$~cm$^3$\,g$^{-1}$\,s$^{-2}$，而 $\alpha^2 = 5.3251\times10^{-5}$ 无量纲。"$G \approx \alpha^2\mu_0$"在 CGS 中\textbf{没有对应表达式}——左边是 $10^{-8}$ 量级的有量纲数值，右边是 $10^{-5}$ 量级的无量纲数值，比较无意义。SI 中的"数值相近"因此是单位制产物（unit artefact）。

\textbf{（二）真正的无量纲不变量}。由 $G$、$\alpha$ 与一个质量（如电子质量 $m_e$）可构造单位制无关的引力耦合常数
\begin{equation}
  \Pi_1 \;=\; \frac{G m_e^2}{\hbar c} \;=\; \frac{m_e^2}{M_P^2} \;=\; 1.7518(5)\times10^{-45},
\end{equation}
即引力的无量纲耦合常数（两电子间引力以 $\hbar c$ 为尺度）。其倒数 $\Pi_1^{-1} = (M_P/m_e)^2 = 5.71\times10^{44}$ 是普朗克质量与电子质量之比的平方；而两电子间\textbf{电磁力与引力之比}还需乘上 $\alpha$：
\[
  \frac{F_{\text{em}}}{F_{\text{grav}}} \;=\; \frac{e^2/(4\pi\varepsilon_0)}{G m_e^2} \;=\; \frac{\alpha}{\Pi_1} \;=\; 4.17\times10^{42},
\]
常见文献给出的 $4.2\times10^{42}$ 正是此值，\textbf{不是} $\Pi_1^{-1}$。$\Pi_1$ 才是可跨单位制讨论"引力与电磁学联系"的合法对象。

\textbf{（三）$\Pi_1$ 与 $\alpha$ 之间不存在\emph{低复杂度}接近 1 的关系}。由 (四) 的机械化结论，本节量纲格中一切无量纲单项式都由 $\Pi_1$ 与 $\alpha$ 生成，所以"可能的巧合"只能形如 $\Pi_1^{p}\alpha^{q}\approx 1$（$p,q\in\mathbb{Z}$）。对 $|p|\le 3,\ |q|\le 3$ 的 $48$ 个非零指数对穷举，离 $1$ 最近的是 $\alpha^{\pm1}$，也差整整 $137$ 倍；区间 $[1/2,2]$ 内一个都没有。（注：(四) 机器证得的两生成元实为 $\Pi_1$ 与 $\mu_0ce^2/\hbar = 4\pi\alpha$，本节用 $\alpha$ 枚举相当于丢掉 $(4\pi)^p$ 因子；本人按 $\Pi_1^{p}(4\pi\alpha)^{q}$ 重新穷举同一范围，$[1/2,2]$ 内仍为 $0$ 个，故结论不受该简化影响。）

\textbf{但必须诚实说明这一穷举口径的边界}。$\ln\Pi_1/\ln\alpha = 20.9452$，其连分数展开的首个收敛点恰为 $21/1$，于是
\[
  \Pi_1\,\alpha^{-21} \;=\; 1.3092 \pm 0.00004,
\]
相对 $1$ 偏 $31\%$，对数间距 $|\ln 1.3092| = 0.269$。须说明这里与 §4.2 的 $0.2616\%$ 不是同一把尺：那里问"\textbf{两个 SI 数值是否重合}"，对数间距 $|\ln(1.002616)| = 0.0026$；这里问"\textbf{是否存在接近 $1$ 的无量纲不变量}"，阈值只能是"多近算近"的人为约定——本文因此不用百分比下断言，而直接给出对数间距，并把它随 $|q|$ 增大而可被压缩这一事实作为不主张全域不可能的理由。更一般地，只要 $\ln\Pi_1/\ln\alpha$ 是无理数（本文未能证明），Dirichlet 逼近定理就保证存在无穷多整数对使 $\Pi_1^{p}\alpha^{q}$ 任意接近 $1$。因此"不存在接近 $1$ 的无量纲不变量"这一类断言\textbf{只有在显式给定复杂度上限后才有意义}，脱离上限的全域不可能陈述是本论文不予主张的（参见 §4.4 对事后选择偏差的同类讨论）。反之，若要通过无量纲化"挽救" $G \approx \alpha^2\mu_0$，必须引入比值 $G/(\alpha^2\mu_0)$ 本身——但有量纲量之比依赖单位制，不是不变量。

\textbf{（四）上述结论的 Lean 形式化}。"不存在无量纲不变量"此前只是解析推导，本文将其升级为机器检查的定理（\texttt{G\_dimension\_lattice.lean}，不依赖 Mathlib，仅用 Lean core 的 \texttt{Int} 与 \texttt{omega} 决策过程）。把量纲实现为 SI 四基本量纲 $(M,L,T,I)$ 上的整数指数向量，对单项式 $G^a \mu_0^{c} c^{d} \hbar^{f} e^{g} m^{h}$ 机器证明：

\begin{enumerate}
  \item \texttt{dimMono\_param}：该单项式无量纲 $\Rightarrow$ $g = 2c \wedge f = -a-c \wedge d = c-a \wedge h = 2a$。即无量纲单项式的指数格是\textbf{秩 2} 的自由交换群，两个生成元恰为 $\Pi_1 = Gm^2/(\hbar c)$（取 $a=1,c=0$）与 $\mu_0 c e^2/\hbar = 4\pi\alpha$（取 $a=0,c=1$）。
  \item \texttt{no\_dimensionless\_G}（\S4.7 主结论）：\textbf{不引入质量标度}（$h=0$）时必有 $a=0$——$G$ 不可能出现在任何无量纲单项式里。这把"找不到单位制无关的 $G$–$\alpha^2\mu_0$ 关系"从经验陈述提升为线性代数上的不可能性。
  \item \texttt{Pi1\_dimensionless} 与 \texttt{alpha\_dimensionless}：$\Pi_1$ 与 $\alpha$ 的无量纲性由同一条判定式机器验证，说明 (二) 中构造合法、且障碍恰是"缺一个质量"。
  \item \texttt{dimG\_ne\_dim\_mu0}：$[G] \neq [\mu_0]$；因 $\alpha$ 无量纲，$[\alpha^2\mu_0] = [\mu_0]$，故 SI 中的量纲矛盾是\textbf{可判定}事实。
\end{enumerate}

\noindent 附录~D 清单中的全部定理均已逐条 \texttt{\#print axioms} 核验（2026-10-16 实测），九条全部无 \texttt{sorryAx}：
\begin{itemize}
  \item \texttt{thm1\_monoid}：\textbf{不依赖任何公理}；
  \item 五条量纲判定链依赖 \texttt{propext}、\texttt{Classical.choice} 与 \texttt{Quot.sound} 三条基础公理：
        \texttt{dimMono\_eq\_zero\_iff}； \texttt{dimMono\_param}； \texttt{Pi1\_dimensionless}；
        \texttt{alpha\_dimensionless}； \texttt{no\_dimensionless\_G}。
  \item \texttt{dimG\_ne\_dim\_mu0}、\texttt{cert\_gap\_pos}、\texttt{alpha2mu0\_gt\_G\_of\_cert}：依赖 \texttt{[propext, Quot.sound]}。
\end{itemize}
\noindent 仍需说明其边界：形式化覆盖的是\textbf{单项式}（幂积）层面的不变量分类，不是任意表达式；把结论推广到含超越函数的表达式类需要额外的量纲不变性假设，本文不做此推广。

\textbf{结论}：\textbf{不存在}在单位制变换下保持不变、且数值上体现"$G$ 与 $\alpha^2\mu_0$ 联系"的无量纲组合（单项式情形已机器证明）。$G \approx \alpha^2\mu_0$（SI）是纯 SI 单位制下的数值巧合，不具有单位制无关的物理意义。若引力与电磁学间存在真实联系，其表述对象应是 $\Pi_1$（或其质子-电子版本 $\Pi_{pe} = Gm_pm_e/(\hbar c) = 3.2\times10^{-42}$）与某个待求理论值之间的关系，而非 $G$ 与 $\mu_0$ 的数值比较。

\subsection{认证区间验证：$0.26\%$ 偏差为严格事实，相当于 $116\sigma$}

启发式浮点容差不足以支撑"不相交"这类否定性结论，必须以认证算术替代。本文以 \textbf{Arb 球算术}\cite{johansson2017}（经 \texttt{python-flint} 0.9.0 绑定\cite{pythonflint}，中点精度 128 bit；随机抽样环节使用 \texttt{numpy}\cite{numpy}）为主验证手段：Arb 的每一步运算都按\textbf{定向舍入}扩展半径，因此输出区间是数学上严格的包含集，而非误差传播的启发式包络。以 CODATA 2022 的中心值 $\pm$ 标准不确定度为输入球（脚本把 $\alpha$ 与 $\mu_0$ 当\textbf{独立}球输入，半径相加，给 $u_r(\alpha^2\mu_0)=4.61\times10^{-10}$，比完全相关的 $4.52\times10^{-10}$ 宽 $1.86\%$。\textbf{这 $1.86\%$ 几乎不是独立性带来的代价}：把 $\mu_0$ 端球换成由 $\alpha$ 传递的$1.89428\times10^{-16}$ 后，独立 trace 给 $4.52221\times10^{-10}$，与完全相关的 $4.52226\times10^{-10}$ 只差 $1.0\times10^{-5}$（相对）——本文实测，工件见 \texttt{results/q5b\_alpha\_ball\_nist\_caliber.json} 的 \texttt{independence\_vs\_display}。所以 $1.86\%$ 里有 $1.85\%$ 来自 CODATA 把 $u(\mu_0)$ 向上舍入展示成 $(20)$，独立性本身只值 $10^{-5}$ 量级。另勿与 §4.2 混淆：那里是\textbf{误差传播公式}误把相关分量当独立做平方和，会\textbf{低估}为 $\sqrt{5}\,u_r(\alpha)$；这里是\textbf{区间算术}，半径相加只会\textbf{高估}宽度。两种形式体系方向相反，但不矛盾）（脚本 \texttt{src/q3\_arb\_certified.py}，输出 \texttt{results/q3\_arb\_results.json}）：
\begin{align}
  \alpha^2\mu_0 &\in \bigl[\,6.6917625530\times10^{-11},\; 6.6917625593\times10^{-11}\,\bigr], \\
  G &\in \bigl[\,6.6741499999\!\times10^{-11},\; 6.6744500001\!\times10^{-11}\,\bigr].
\end{align}
两球\textbf{不相交}（由 Arb 原生 \texttt{overlaps()} 判定为假）。其差的认证区间为
\[
  \alpha^2\mu_0 - G \;\in\; \big[\,1.73125\times10^{-13},\ 1.76126\times10^{-13}\,\big] \;=\; (1.7463 \pm 0.0150)\times10^{-13},
\]
\textbf{下界严格为正}（$1.73125\times10^{-13} > 0$）。\textbf{本节所有展示区间按所列位数向外取整}（下界取 floor、上界取 ceil），故展示区间恒包含认证区间；本文自查中一度写成 $1.73126$ 与 $6.6917625531$，那已是向内取整，现按纪律改正。该下界正是两区间之间的间距，相当于 $G$ 标准不确定度半宽的 $115.4$ 倍。故在 CODATA 2022 的不确定度范围内，$\alpha^2\mu_0 > G$ 是认证可证的严格事实。相对偏差的认证区间为
\[
  \frac{\alpha^2\mu_0 - G}{G} \in \bigl[\,0.2593\%,\; 0.2639\,\%\bigr] = 0.2616\% \pm 0.0023\%,
\]
半宽 $0.0023\%$ 完全来自 $G$ 端（$\alpha^2\mu_0$ 一侧仅贡献 $\sim 5\times10^{-8}\%$）。

\textbf{完全相关的第二份证书}。由于 SI 中 $\mu_0 = 2\alpha h/(e^2c)$ 且 $h,e,c$ 为精确值，可在\textbf{同一条算术 trace} 内由 $\alpha$ 直接导出 $\mu_0$，于是 $\alpha^2\mu_0 = 2\alpha^3 h/(e^2c)$，其不确定度\textbf{只来自 $\alpha$}（相对 $3u_r(\alpha)$）。机器实测（脚本 \texttt{src/q5\_correlated\_and\_spacing.py}，结果 \texttt{results/q5\_correlated\_and\_spacing.json}）：$\alpha^2\mu_0 \in [6.6917625531, 6.6917625592]\times10^{-11}$（$u_r(\alpha^2\mu_0) = 4.522\times10^{-10}$），与 $G$ 的区间\textbf{仍不相交}，间距 $1.731255\times10^{-13} = 115.42\,u(G)$，偏差 $0.261639\%$（$116.4\sigma$）；由 $\alpha$ 导出的 $\mu_0$ 中心值为 $1.256\,637\,061\,2628\times10^{-6}$，与 CODATA 展示值 $1.256\,637\,061\,27(20)\times10^{-6}$ 的相对差为 $5.7\times10^{-12}$。这个差\textbf{单靠 $\mu_0$ 自身的展示舍入解释不了}（其展示末位半 ulp 为 $0.5\times10^{-17}$，对应相对 $4.0\times10^{-12}$，实测差是它的 $1.44$ 倍），需再动用“$\alpha$ 中心值也只展示到 10 位”（半 ulp 对应相对 $6.9\times10^{-12}$）才能容纳。更要说清的是：该差只有 $u(\mu_0)=2.0\times10^{-16}$ 的 $1/28$，\textbf{完全落在 $\mu_0$ 自己的误差带内}，所以它只是“推导链没有被展示位数证伪”的弱核对，不是对精度或物理自洽性的一般性检验。两版证书的\textbf{结论完全一致}，差别仅在 $\alpha^2\mu_0$ 一侧的区间宽度（$4.52$ vs $4.61\times10^{-10}$，归因见本节开头那条括号注：$1.86\%$ 里 $1.85\%$ 是 CODATA 对 $u(\mu_0)$ 的向上舍入展示，独立性只值 $0.006\%$）。由于该侧比 $G$ 端窄 $4.97\times10^{4}$ 倍（接近五个数量级），这一宽度差对“是否重合”的判定不起决定作用。\par\medskip\noindent\textbf{该认证陈述的确切范围}：本节认证的是\textbf{在 SI 数值口径下}"$\alpha^2\mu_0$ 的数值不等于 $G$ 的数值"。它被排除的候选命题只有一种——"$G$ 的计量不确定度尚未收敛，测精之后两数会重合"。它\textbf{不}给出任何单位制无关的物理差异：§4.7 已表明"$G$ 与 $\alpha^2\mu_0$ 数值相近"这一说法本身依赖单位制，故 $116\sigma$ 是 SI 口径内的严格性，而不是"发现了一种新的引力--电磁联系（或其缺失）"的证据。真正的否定性结论由 §4.7 的量纲不可比性承担，§4.8 只负责把数值层面的"可能重合"也一并关掉。\textbf{注}：python-flint 的 \texttt{str(arb)} 显示精度低于端点实值，本文报告的一切区间均取自 \texttt{lower()/upper()}。\par\medskip\noindent\textbf{两种口径必须分开说}（本节的"不相交"与"σ 水平"用的是\textbf{不同的}不确定度语义，不写明就会自相矛盾）：其一，\textbf{认证口径}把 $\mathrm{mid}\pm u$ 当作\textbf{硬夹逼}（即真值必在该球内），此时"两区间不相交 $\Rightarrow \alpha^2\mu_0 > G$"是区间算术内的严格推理；其二，\textbf{统计口径}把 $u$ 当作 $1\sigma$ 概率带，此时才有"$116\sigma$""$1.7\sigma$ 与相等相容"这类说法。本文两者都用，但\textbf{不把统计口径的显著性等价于认证口径的定理性}；若采用统计口径，则结论强度取决于 $u(G)$ 的取值（见下文稳健性段）。\par\medskip\noindent\textbf{$\sigma$ 水平}：以 $G$ 的标准不确定度 $u^{(2022)}(G)=1.5\times10^{-15}$（即 $2.247\times10^{-5}$ 的相对值）为单位，偏差 $2.6164\times10^{-3}$ 对应
\begin{equation}
  \frac{|\alpha^2\mu_0 - G|}{u(G)} \;=\; \frac{2.6164\times10^{-3}}{2.247\times10^{-5}} \;\approx\; \mathbf{116\sigma}.
\end{equation}
（若改用 CODATA 2022 展示的舍入值 $u_r^{(2022)}(G)=2.2\times10^{-5}$，得 $119\sigma$；两者均远超任何显著性阈值。本式左端的 $2.247\times10^{-5}$ 是\textbf{未舍入}值 $u/G$，二者不矛盾，只是位数不同。）

\textbf{对 $G$ 未来调整的稳健性}：$G$ 是唯一"多组高精度实验彼此不相容"的基本常数。过去二十年间的 11 个精密测量值被 CODATA 2014 全部收录，推荐值 $G = 6.67408(31)\times10^{-11}$ 的相对不确定度 $u_r^{(2014)}(G) = 4.7\times10^{-5}$，是各基本常数中最大者之一\cite{li2020g}；更早的 CODATA 1998 调整因各次测量互相矛盾，把 $u_r^{(1998)}(G)$ 从 1986 年的 $1.28\times10^{-4}$ \textbf{放大}到 $1.5\times10^{-3}$（约 $12$ 倍）\cite{codata1998}。结论必须对这种膨胀做压力测试：

其一，\textbf{中心值层面}：偏差 $2616\,\mathrm{ppm}$（$=2.6164\times10^{-3}$）远大于历次 CODATA 推荐值之间的移动——2014 $\to$ 2022 仅 $33\,\mathrm{ppm}$，而历史上最大的一次（1998 $\to$ 2002，$6.673(10)\to 6.6742(10)\times10^{-11}$）也只有约 $180\,\mathrm{ppm}$；两者都小于偏差的 $1/14$，故推荐值本身如何移动都不会与 $\alpha^2\mu_0$ 相遇。其二，\textbf{不确定度层面}：若把 $u(G)$ 人为取到史上最夸张的 $1500\,\mathrm{ppm}$（1998 年情形），显著性由 $116\sigma$ 降为 $2616/1500 \approx 1.7\sigma$。必须诚实承认：$1.7\sigma$（双侧 $p \approx 0.09$）\textbf{落在"与相等假设相容"的范围内}，那种情形下"$\alpha^2\mu_0 \neq G$"将不再是一条可断言的事实。可见 $116\sigma$ 是\textbf{依赖计量能力}的结论，其成立前提是当前 $G$ 的相对不确定度仅为 $2.2\times10^{-5}$，它不是先验定理、也不免疫未来测量格局的变化。但即使真退到 $1.7\sigma$，§4.7 的量纲不可比性\textbf{不受影响}——数值上的相容也不会使"$G = \alpha^2\mu_0$"成为单位制无关的陈述。

\textbf{与 mpmath 包络的一致性}：作为交叉核对，\texttt{mpi}（误差传播包络，非逐舍入认证）给出同样的不相交结论，间距 $1.7312553097\times10^{-13}$ 与 Arb 的 $1.7312553097\times10^{-13}$ 在小数点后 8 位一致。\textbf{一致的前提是两者用同一组输入球}：本文一度让 R4 用 $u_r(\alpha)=1.6\times10^{-10}$（即 $u(\alpha)=1.1676\times10^{-12}$），那时两法只在第 6 位一致，残差 $1.2\times10^{-21}$ 纯属口径差而非算术差；现已把 R4 与 q3/q5 对齐为展示末位口径（脚本 \texttt{R4\_interval\_arithmetic.py}，工件 \texttt{results/mpi\_crosscheck.json}）。

\textbf{同一套 Arb 证书给出的其他严格判定}：$\Lambda - \alpha^{-1} = (3.0460 \pm 0.0207)\times10^{-4}$（端点区间 $[3.04578, 3.04619]\times10^{-4}$）\textbf{不含} $0$（故 $\Lambda = \alpha^{-1}$ 为假，不只是"近似"）；$\alpha^2\Lambda^2 = 1.000004446 \pm 7.8\times10^{-10}$ \textbf{不含} $1$；$M^*c^2 \in [8.9091, 8.9093]\times10^{16}$~GeV。

\textbf{证书回放进 Lean}：区间端点由 Arb 在 Lean 之外算出，但"给定证书 $\Rightarrow$ 结论"这一步已机械化——把 $\alpha^2\mu_0$ 的下界与 $G$ 的上界以整数形式（单位 $10^{-22}$，向外取整以保证包含性）写入 \texttt{G\_dimension\_lattice.lean}，机器验证 \texttt{cert\_gap\_pos}（间距 $>0$）与回放引理 \texttt{alpha2mu0\_gt\_G\_of\_cert}（$\forall x\,y,\ x \ge \text{下界} \to y \le \text{上界} \to x > y$）。于是留在 Lean 之外的只有"端点来历"一项，而该项由 Arb 的定向舍入语义负责。

\textbf{仍未完成的部分}：其一，上述区间是\textbf{输入不确定度}与\textbf{算术舍入}的联合包络，它认证的是"给定 CODATA 输入下的推导"，而非把 CODATA 调整本身重新证明一遍；其二，Arb 侧的完整依赖图尚未导出为机器可读证书，Lean 内也未重放 Arb 的运算步骤（只重放了单调比较），Lean 生态的认证浮点库 FloatLib 仍不成熟。

%==============================================================================
\section{讨论}
%==============================================================================

\subsection{验证结果的综合评估}

表~\ref{tab:summary} 汇总了原文六项核心断言的验证结果，以及本文新增了无量纲不变量分析（§4.7）与认证区间检验（§4.8）两项补充检验。

\begin{table}[h]
\centering
\caption{验证结果汇总}
\label{tab:summary}
\footnotesize
\setlength{\tabcolsep}{2pt}
\begin{tabular}{@{}llcl@{}}
\toprule
\textbf{断言} & \textbf{内容} & \textbf{验证方式} & \textbf{性质} \\
\midrule
定理 1 & $\alpha^2/(c^2\varepsilon_0) = \alpha^2\mu_0$ & Lean 形式证明 & 代数恒等式（平凡） \\
定理 2 & $G \approx \alpha^2\mu_0$ (SI) & Lean + mpmath + Arb & 单位制依赖的数值巧合 \\
定理 3 & $[G] = [M]^{-2}$, $[\alpha^2\mu_0]_{\text{nat}} = 1$ & 解析 + Lean & 量纲分析（需重释 $\mu_0$） \\
定理 4 & $\Lambda \approx 137 \approx \alpha^{-1}$ & 穷举 + 零模型 & 全族最优；$p$ 非内禀（§4.4） \\
定理 5 & $M^*c^2 \approx E_\text{GUT}$ & Lean + mpmath & 定义导出；靶值模型依赖 \\
定理 6 & $G \approx \alpha^2/M^{*2}$ & Lean + mpmath & 由 $\Lambda \approx \alpha^{-1}$ 保证 \\
§4.7 & $\Pi_1 = Gm_e^2/\hbar c = 1.7518(5)\times10^{-45}$ & Lean + Arb & 唯一合法不变量，与 $\alpha$ 不贴近 \\
§4.8 & $\alpha^2\mu_0$ 与 $G$ 区间不相交 & Arb 球算术 & SI 内不相交，$116\sigma$ \\
\bottomrule
\end{tabular}
\end{table}

\subsection{关于"量纲矛盾的解决"}

原文的核心主张是 $G = \alpha^2\mu_0$ 的"量纲矛盾"可以通过切换到自然单位制来解决。本文的量纲分析（定理~3）给出了比原文更精确的结论：在标准自然单位制（Heaviside--Lorentz 约定）下，$[\mu_0] = 1$（无量纲），而 $[G] = [M]^{-2}$，因此 $G = \alpha^2\mu_0$ 在自然单位制中\textbf{仍然不量纲齐次}。量纲矛盾并未因切换单位制而消失，而是从 SI 中的两种复合量纲不等转化为自然单位制中的 $[M]^{-2} \neq 1$。

要使 $G \approx \alpha^2\mu_0$ 成为量纲齐次的关系式，需要\textbf{重新诠释} $\mu_0$ 的量纲为 $[M]^{-2}$，而非采用标准电磁学中的无量纲 $\mu_0$。这是一种物理假设（将 $\mu_0$ 视为与引力耦合相关的导出量），而非单纯切换单位制即可自动实现的效果。

\textbf{标准场论立场}：在标准量子场论教材（如 Peskin \& Schroeder\cite{peskin1995}）的 Heaviside--Lorentz 自然单位制中，$\mu_0$ 与 $\varepsilon_0$ 被\textbf{完全消除}——它们不出现于拉格朗日量，电磁耦合的全部信息集中在无量纲的 $\alpha$ 中。因此"给 $\mu_0$ 赋 $[M]^{-2}$ 量纲"等价于重新引入一个已被约化的量作为独立标度场，需要具体动力学模型支撑；本文不提出此类模型，故 $[\mu_0] = [M]^{-2}$ 只能作为\textbf{形式上的量纲记账}，不构成物理机制。（SI 基本单位与导出单位的区分在计量学界仍有公开讨论，见 \cite{brown2023}。）

\textbf{思想史坐标}：值得指出的是，"引力常数是电磁真空性质的宏观残留"并非无传统的 ad hoc 假设。Sakharov (1967)\cite{sakharov1967}的诱导引力纲领提出 $G^{-1}$ 由量子场真空涨落的二阶项决定，把引力视为电磁（及更一般量子场）真空的次生现象——这是"$G$ 可由 $\alpha$、$\mu_0$ 等电磁量导出"一类想法在物理学界最严肃的先驱。其思想传承可上溯普朗克（自然单位制的必要性）、Kaluza--Klein（电磁与引力的几何统一）、惠勒（质量与电荷的几何起源）、爱丁顿（$\alpha$ 的拓扑起源猜想），形成一条"引力非基本"的假说谱系。\textbf{但这一定位是动机性的，不是论证性的}：Sakharov 机制给出的是 $G^{-1} \sim E_{\text{UV}}^2$ 型量子场能标关系（$E_{\text{UV}}$ 为紫外截断能标，与定义~1 的几何常数 $\Lambda$ 无关），并不预言 $G = \alpha^2\mu_0$ 的具体数值形式；从"有思想传统"到"该公式成立"之间仍差一个可检验的动力学模型。要使 $G \approx \alpha^2\mu_0$ 物理化，可能的理论结构包括：额外维度/Kaluza--Klein 型统一、标量-张量引力中 $G$ 作为动力学场、以及某种紫外-红外对应。本文仅指出这些是必要条件，不探索其中任何一个。

\subsection{$G \approx \alpha^2\mu_0$ 在物理学常数关系研究中的地位}

物理常数之间的数值巧合在物理学史上并不罕见。Eddington（1930 年代）\cite{eddington1936}试图从纯组合论证 $\alpha^{-1}$ 必为\textbf{整数}——先取 $136$，测量精度提高后又"加一"改为 $137$，此事即其学生戏称其 "Arthur Adding-One" 的由来；他并构建了电磁学与引力的比率 $\alpha/(Gm_pm_e/\hbar c) \sim 2\times10^{39}$（质子-电子）等无量纲大数。Dirac (1937) 的大数假说\cite{dirac1937}则假设这些大数随宇宙年龄变化，引发了关于常数是否真正"恒定"的持续讨论，其当代系统综述见 \cite{uzan2011}。同一类构造在当代民间物理写作中仍在延续，例如与本文原文同期出现的"四大常数关系"方案\cite{zhangxiangqian}——按本文同样的量纲分析，它同样不齐次、同样依赖单位制，且其"$m = qL/T$"一类的量纲赋值属循环论证。这类工作的共同问题不是"想法大胆"，而是\textbf{缺少本文 §4.7--4.8 的两道闸门}：先要有一个单位制无关的无量纲对象，再要有一个认证区间说明它与实验值是否真的可区分。

本文构造的 $\Lambda = 4\pi^3+\pi^2+\pi \approx \alpha^{-1}$ 以及由此导出的 $G \approx \alpha^2\mu_0$（SI 中偏差 $0.26\%$），在性质上与 Eddington 的 $\alpha^{-1} \approx 137$ 猜测同类：它们都是基于纯数学构造（$\pi$ 的多项式）对实验测量值的数值逼近。§4.4 的穷举加零模型检验进一步表明：虽然 $\Lambda$ 是同族 $194\,480$ 个候选中的全局最优，但在"靶值与格点无关"的零假设下出现同量级最优拟合的概率 $p \approx 0.13$（最保守口径 $0.34$）——在该家族口径下这类吻合\textbf{统计不显著}。同时如实报告其局限：把家族收紧到 $\lvert a_i\rvert\le 4$ 时同一算法给出 $p\approx 0.011$，即"$\Lambda$ 命中"\textbf{在此界下反而是小概率事件}；两种读法都只在给定家族与给定零模型内成立，本文因此不把任何一个 $p$ 当作结论性证据（见 §4.4 缩族段）。诚实的定性是"数值巧合"：其价值在于启发思考，但不能替代物理解释。同一方法论教训也见于音乐律学中的毕达哥拉斯逗号研究——$(3/2)^{12}/2^7 = 531441/524288 \approx 1.0136$ 一类"简单整数比逼近"同样常被误读为深层秩序，最终须以"承认巧合"处理。

值得强调的是（§4.8 的区间证明）：$0.26\%$ 的偏差远大于 $G$ 的测量不确定度（$2.2\times10^{-5}$），两常数的不确定区间严格不相交，因此这一偏差是\textbf{真实的数值差异}而非测量误差所致。必须限定：该说法成立于\textbf{当前} $u_r(G)=2.2\times10^{-5}$ 的测量水平；若 $G$ 的不确定度退回 1998 年式的 $1.5\times10^{-3}$，两者在约 $1.7\sigma$ 内反而相容（§4.8）——那是测量能力的问题，不是"两数将重合"的证据。

\subsection{关于 $G = \alpha^2/M^{*2}$ 的地位}

定理~6 表明 $G \approx \alpha^2/M^{*2}$ 在自然单位制中高度自洽（偏差仅 $4.45\times10^{-6}$）。但这一自洽性的来源是：

\begin{enumerate}
  \item $M^*$ 由 $M^* = M_P/\Lambda$ 定义；
  \item $\Lambda \approx \alpha^{-1}$（定理~4 的数值巧合）；
  \item 因此 $\alpha^2/M^{*2} = \alpha^2\Lambda^2/M_P^2$，而自然单位制下 $G_{\rm nat} \equiv 1/M_P^2$ \textbf{是定义}。于是 $\alpha^2/M^{*2} \approx G$ 与 $\alpha^2\Lambda^2 \approx 1$ 是\textbf{同一个命题的两种写法}——定理~6 声称的"偏差仅 $4\times10^{-6}$"与定理~4 的 $|\alpha^2\Lambda^2 - 1| = 4.4\times10^{-6}$ 是\textbf{同一个数}，不是两次独立验证。
\end{enumerate}

换言之，$G \approx \alpha^2/M^{*2}$ 的成立完全依赖 $\Lambda \approx \alpha^{-1}$：把 $M^* \equiv M_P/\Lambda$ 代入即得$\alpha^2/M^{*2} = \alpha^2\Lambda^2 \cdot G_{\rm nat}$，两者之差是恒等的代数重排，$G/(\alpha^2/M^{*2}) = 1/(\alpha^2\Lambda^2)$ 也随之一并确定。所以这条"自洽性"链条不含任何独立信息量：它不会因物理测量而变，只因 $\Lambda$ 的定义而变；把它列为独立证据即构成 §4.1 所批评的"同一恒等式重复计数"。本文据此把定理~6 判定为\textbf{代数推论}而非物理定律。

\subsection{关于 $M^*c^2$ 与 GUT 能标的对应}

定理~5 中 $M^*c^2 = 8.909\times10^{16}$~GeV 的数值本身无争议，但"对应大统一能标"这一表述\textbf{依赖靶值的选取}。GUT 能标不是自然界的单一可观测量，而是具体统一模型（含其粒子含量与重整化群跑动）的拟合输出：最小非超对称 SU(5) 的规范耦合统一能标约 $10^{15}$~GeV，最小 SO(10) 约 $10^{16}$~GeV，而 MSSM/SUSY SU(5) 情形约 $2\times10^{16}$~GeV（表~\ref{tab:gut}）。$8.9\times10^{16}$~GeV 落在所有主流模型能标的同数量级区间（$10^{15}$--$10^{17}$~GeV）偏上端，但与任何一个具体模型都不吻合到 $1\%$——原文的 $<1\%$ 偏差是相对"$8.9\times10^{16}$"这一未指明模型出处的靶值计算的。若不指定模型，"GUT 能标"只是一个跨三个数量级的范围，$M^*c^2$ 落入其中是平凡结论。

此外，这一对应链条本身依赖于 $\Lambda \approx \alpha^{-1}$：
\[
  M^*c^2 = \frac{M_P c^2}{\Lambda} \approx \frac{M_P c^2}{\alpha^{-1}} = \alpha \cdot M_P c^2 \approx \frac{1}{137} \times 1.22\times10^{19}~\text{GeV} \approx 8.9\times10^{16}~\text{GeV}.
\]
这一数值链条的起点是 $\alpha \approx 1/137$，终点是 GUT 能标量级，中间经过纯几何常数 $\Lambda$ 的桥梁作用。换言之，$M^*c^2$ 与 GUT 能标量级的"对应"归根结底是 $\alpha^{-1} \approx 137$ 与 $M_P c^2 \approx 1.22\times10^{19}$~GeV 这两个已知事实的商，不含独立的物理内容。

%==============================================================================
\section{结论}
%==============================================================================

本文使用 Lean 4 定理证明器、Python mpmath 高精度计算（CODATA 2022 常数）、Arb 定向舍入认证区间与全族穷举加零模型检验，对塔里芙 (2026) 关于 $G \approx \alpha^2\mu_0$ 的全部核心断言进行了多层次验证。主要结论如下：

\begin{enumerate}
  \item \textbf{代数恒等式（平凡）}：$\alpha^2/(c^2\varepsilon_0) = \alpha^2\mu_0$ 由 $c^2 = 1/(\mu_0\varepsilon_0)$ 直接保证，不涉及物理假设。原文将此式与 $G = \alpha^2\mu_0$ 写成两段，制造了"双重验证"的错觉。

  \item \textbf{数值巧合（单位制依赖）}：$G \approx \alpha^2\mu_0$（CODATA 2022 下偏差 $0.2616\%$）在 SI 中成立，但为单位制依赖的数值关系。CGS 高斯制中该关系无对应表达式（表~\ref{tab:units}）；在单项式类内不存在承载该联系的无量纲不变量（§4.7，该否定结论已机器证明，推广到任意表达式类未证）；Arb 认证区间证明两不确定区间不相交，偏差相当于 $116\sigma$，是严格事实而非测量误差（§4.8）。

  \item \textbf{量纲分析（需重新诠释，核心部分已机械化）}：在标准自然单位制下，$[G] = [M]^{-2}$ 而 $[\alpha^2\mu_0] = 1$，两者并不相等。要使 $G \approx \alpha^2\mu_0$ 量纲齐次，须重新赋予 $\mu_0$ 以 $[M]^{-2}$ 量纲——这是标准场论中没有依据的物理假设（Sakharov 诱导引力\cite{sakharov1967}提供动机传统但不提供该公式），而非单纯切换单位制的效果。"不引入质量标度则 $G$ 不可能出现在任何无量纲单项式中"已在 Lean 中形式化（\texttt{no\_dimensionless\_G}，附录~D），故 §4.7 的否定结论不再是解析推导而是机器检查的定理。

  \item \textbf{几何-经验吻合（显著性随家族边界双向漂移，不是内禀量）}：$\Lambda = 4\pi^3+\pi^2+\pi = 137.03630$ 与 $\alpha^{-1} = 137.035999$ 仅差 $3.0\times10^{-4}$，是全部数值关系的枢纽。$194\,480$ 个低复杂度 $\pi$ 多项式全族穷举显示 $\Lambda$ 为同族全局最优（次优差 $12$ 倍），但在显式零模型下 $p \approx 0.13$（各口径 $0.123$--$0.340$ 一致），\textbf{统计不显著}；反向收紧家族（$\lvert a_i\rvert\le4$，$N=6\,560$）时同一算法给出 $p\approx 0.011$，此时"$\Lambda$ 恰好命中"\textbf{反而显著}；因此本文对 (d) 的表述口径是"显著性对人为家族边界敏感（$0.011$--$0.73$），不构成内禀证据"，而非单点"$p=0.13$ 不显著"；把基扩到 $(\pi, e)$ 后出现 $-1-2\pi+5\pi^2+3\pi^3-2e+e^2 = 137.036159$，比 $\Lambda$ 好 $1.9$ 倍——"最优"本身是家族规模的函数。Arb 认证区间另给出 $\Lambda \neq \alpha^{-1}$ 的严格判定。该吻合在性质与命运上与 Eddington 对 $\alpha^{-1} \approx 137$ 的纯理论猜测同类。

  \item \textbf{GUT 对应（模型依赖的代数推论）}：$M^*c^2 = 8.909\times10^{16}$~GeV 与主流 GUT 模型能标（最小 SU(5) $\sim 10^{15}$、SO(10) $\sim 10^{16}$、SUSY SU(5) $\sim 2\times10^{16}$~GeV）仅同数量级而不吻合；"<1\%" 偏差是相对未指明模型出处的靶值 $\alpha M_P c^2$ 计算的，为 $\Lambda \approx \alpha^{-1}$ 的直接推论，不构成独立的物理预言。

  \item \textbf{核心公式（代数推论）}：$G \approx \alpha^2/M^{*2}$ 在自然单位制下自洽（$\alpha^2\Lambda^2 = 1.000004$），是 $\Lambda \approx \alpha^{-1}$ 的代数推论，而非独立的物理定律。
\end{enumerate}

\subsection*{局限性与未来工作}

本文存在以下局限：(i)~Lean 形式化覆盖三类内容——代数恒等式（定理~1）、量纲格的单项式不变量分类（§4.7(四)）、Arb 证书回放（§4.8）；\textbf{未}形式化的是 $\mu_0$ 量纲重诠释的物理可接受性（属物理论证）与含超越函数的任意表达式类上的不变量分类；(ii)~本文未讨论原文中可能涉及的波动方程等其他断言；(iii)~$0.26\%$ 偏差的物理来源（若存在）尚未给出解释，本文只确证它是真实差异而非测量误差；(iv)~Arb 证书以整数端点形式进入 Lean，未导出完整依赖图，也未在 Lean 内重放 Arb 的运算步骤（只重放了单调比较）；(v)~GUT 模型能标（表~\ref{tab:gut}）取文献量级估计而非逐模型重整化群拟合值；(vi)~验证工件目前以本地目录交付，尚无版本化公开仓库与持久标识符（DOI）。

未来工作：(a)~把 §4.7 的不变量分类从单项式推广到更宽的表达式类，需在 Lean 中形式化"量纲在函数复合下的不变性"；(b)~导出 Arb 的区间依赖图并在 Lean 内重放单调运算，使 §4.8 成为端到端机器检查证明（Lean 生态的认证浮点库 FloatLib 仍不成熟）；(c)~以 MDL（Minimum Description Length，最小描述长度）或通用编码先验替代"指定家族 + 随机靶"的显著性框架，使复杂度惩罚不依赖人为的家族边界（§4.4 的家族稳健性实验已显示结论对此不敏感）；(d)~对 §4.4 指出的"置换检验在事后挑选型问题上退化"作一般化刻画，明确何种零模型适用于何种构造类；(e)~将本文的"全族穷举 $\to$ 显式零模型 $\to$ 家族稳健性"流程应用于其他常数数值巧合（如质子-电子质量比、毕达哥拉斯逗号类构造）。

源代码与验证工件随本文工件包交付，逐项可重跑。

%==============================================================================
% 参考文献
%==============================================================================

\begin{thebibliography}{99}

\bibitem{talihua2026}
塔里芙.
\textit{$G=\alpha^2\mu_0$量纲与数值矛盾的完整解决}.
微信公众号, 2026.
\url{https://mp.weixin.qq.com/s/chsyXx8N8oseJTaPysNuQA}
\textbf{可及性说明}：2026-10-16 复测该链接返回"该内容已被发布者删除"，本文核对所用文本获取于 2026 年 9 月；下文转述的原文断言以该时点可见文本为准，读者可能已无法复核原文。

\bibitem{lean4}
The Lean Community.
\textit{Lean 4 Theorem Prover}.
\url{https://leanprover.github.io/}

\bibitem{codata2022}
P. J. Mohr, D. B. Newell, B. N. Taylor, E. Tiesinga.
\textit{CODATA Recommended Values of the Fundamental Physical Constants: 2022}.
Rev. Mod. Phys. \textbf{97}, 025002 (2025).
\url{https://doi.org/10.1103/RevModPhys.97.025002}; arXiv:2409.03787;
\url{https://physics.nist.gov/cuu/Constants/}

\bibitem{zhangxiangqian}
张祥前.
\textit{解决四大常数关系量纲匹配问题的一个方案}.
微信公众号, 2026.
\url{https://mp.weixin.qq.com/s/euSSkxbNBSVf2riWxuv2kQ}
（访问日期 2026-10-16；该链接需浏览器通过微信验证页方可读取。）

\bibitem{eddington1936}
A. S. Eddington.
\textit{The Relativity Theory of Protons and Electrons}.
Cambridge: Cambridge University Press, 1936.
后续体系见 \textit{Fundamental Theory}（E. T. Whittaker 编，Cambridge University Press, 1946，身后出版）。

\bibitem{dirac1937}
P. A. M. Dirac.
\textit{The Cosmological Constants}.
Nature, 139 (1937): 323.
\url{https://doi.org/10.1038/139323a0}

\bibitem{sakharov1967}
A. D. Sakharov.
\textit{Vacuum Quantum Fluctuations in Curved Space and the Theory of Gravitation}.
Dokl. Akad. Nauk SSSR \textbf{170} (1966): 70；英译 Sov. Phys. Doklady \textbf{12} (1967): 1040；
重印于 Gen. Relativ. Gravit. \textbf{32} (2000): 365--367.
\url{https://doi.org/10.1023/A:1001947813563}
（诱导引力：引力作为量子真空涨落的次生效应。）

\bibitem{peskin1995}
M. E. Peskin, D. V. Schroeder.
\textit{An Introduction to Quantum Field Theory}.
Westview Press, 1995.（本文仅取其 Heaviside--Lorentz 自然单位制约定；不主张具体章节编号。）

\bibitem{johansson2017}
F. Johansson.
\textit{Arb: Efficient Arbitrary-Precision Midpoint-Radius Interval Arithmetic}.
IEEE Transactions on Computers, 66(8) (2017): 1281--1292.
\url{https://doi.org/10.1109/TC.2017.2690633}
（中点--半径球算术与定向舍入；本文 §4.8 经 \texttt{python-flint} 0.9.0 调用其 \texttt{arb} 类型。）

\bibitem{li2020g}
C. Xue, J.-P. Liu, Q. Li, J.-F. Wu, S.-Q. Yang, Q. Liu, C.-G. Shao, L.-C. Tu, Z.-K. Hu, J. Luo.
\textit{Precision Measurement of the Newtonian Gravitational Constant}.
National Science Review, 7 (2020): 1803--1817.
\url{https://doi.org/10.1093/nsr/nwaa165}
（摘要可稽：过去二十年 11 个精密测量值、CODATA 2014 推荐值 $6.67408(31)\times10^{-11}$、$u_r = 47\,\mathrm{ppm}$。）

\bibitem{codata1998}
P. J. Mohr, B. N. Taylor.
\textit{CODATA Recommended Values of the Fundamental Physical Constants: 1998}.
J. Phys. Chem. Ref. Data \textbf{28}(6) (1999): 1713--1852.
同一调整的期刊版见 Rev. Mod. Phys. \textbf{72} (2000): 351--495。
\url{https://doi.org/10.1063/1.556049}
（其中 $G$ 项相对不确定度由 1986 年的 $128\,\mathrm{ppm}$ 放大为 $1500\,\mathrm{ppm}$，以吸收各次测量间的不相容性。）

\bibitem{nath2007}
P. Nath (Pran Nath), P. Fileviez P\'erez (Pavel Fileviez P\'erez).
\textit{Proton Stability in Grand Unified Theories, in Strings and in Branes}.
Physics Reports \textbf{441}(5--6) (2007): 191--317.
\url{https://doi.org/10.1016/j.physrep.2007.02.010}

\bibitem{brown2023}
R. J. C. Brown (Richard J. C. Brown).
\textit{On the Distinction between SI Base Units and SI Derived Units}.
Metrologia \textbf{61}(1) (2024): 013001（在线发表 2023-12-19，印刷卷 2024；本文键名记作 brown2023 取在线日期）。
\url{https://doi.org/10.1088/1681-7575/ad141b}

\bibitem{pythonflint}
F. Johansson \textit{et al.} \texttt{python-flint}: Python bindings for FLINT and Arb, version 0.9.0.
\url{https://github.com/flintlib/python-flint}

\bibitem{numpy}
C. R. Harris, K. J. Millman, S. J. van der Walt, R. Gommers, P. Virtanen, D. Cournapeau \textit{et al.}
\textit{Array Programming with NumPy}.
Nature \textbf{585} (2020): 357--362. \url{https://doi.org/10.1038/s41586-020-2649-2}（本文使用 numpy 2.x 数值库）

\bibitem{mpmath}
The mpmath development team.
\textit{mpmath: a Python library for arbitrary-precision floating-point arithmetic}.
\url{https://mpmath.org/}

\bibitem{uzan2011}
J.-P. Uzan.
\textit{Varying Constants, Gravitation and Cosmology}.
Living Reviews in Relativity \textbf{14}: 2 (2011).
\url{https://doi.org/10.12942/lrr-2011-2}

\end{thebibliography}

%==============================================================================
% 附录：Lean 4 验证代码摘要
%==============================================================================

\appendix
\section{Lean 4 验证代码摘要}

\begin{lstlisting}[language={},caption={物理常数定义（CODATA 2022）}]
def pi_f    : Float := 3.14159265358979323846
def c_si    : Float := 299792458.0          -- 精确
def h_si    : Float := 6.62607015e-34       -- 精确
def mu0_si  : Float := 1.25663706127e-6     -- CODATA 2022, u_r = u_r(alpha) = 1.5e-10
def eps0_si : Float := 8.85418781879e-12    -- CODATA 2022 导出值
def alpha_f : Float := 7.2973525643e-3      -- CODATA 2022
def G_si    : Float := 6.67430e-11          -- u_r ~ 2.2e-5
def hbar_si : Float := h_si / (2.0 * pi_f)
def M_P_si  : Float := Float.sqrt (hbar_si * c_si / G_si)
def Lambda_val : Float := 4.0 * pi_f^3 + pi_f^2 + pi_f
def M_star_si  : Float := M_P_si / Lambda_val
\end{lstlisting}

\begin{lstlisting}[language={},caption={核心验证断言}]
-- 定理2: SI 中 G ~ alpha^2 mu_0, 偏差 0.2616% (CODATA 2022)
def si_dev_pct := Float.abs (G_si - alpha_f^2 * mu0_si) / G_si * 100.0
-- 验证: si_dev_pct = 0.2616 < 0.27  PASS

-- 定理4: Lambda ~ 137 ~ alpha^{-1}
-- 验证: |Lambda - 137| = 0.036 < 0.1  PASS
-- 验证: |Lambda - alpha^{-1}| = 0.0003 < 0.01  PASS

-- 定理5: M*c^2 ~ GUT
-- 验证: M*c^2 = 8.909e16 GeV, |...-8.9e16|/8.9e16 < 0.02  PASS

-- 定理6: alpha^2 Lambda^2 ~ 1
-- 验证: alpha^2 Lambda^2 = 1.000004, |...-1| < 0.002  PASS
\end{lstlisting}

\section{修订新增验证脚本（节选）}

以下脚本位于仓库 \texttt{src/} 目录、完整输出位于 \texttt{results/} 目录，可直接运行复现（依赖 Python 3.13 + mpmath + numpy + python-flint）。

\begin{lstlisting}[language=Python,caption={Lambda 全族穷举（R3 / 定理 4 统计显著性）}]
# R3_lambda_search.py (节选; 完整见 src/)
from mpmath import mp, mpf, pi, fabs
mp.dps = 50
ALPHA = mpf("7.2973525643e-3")      # CODATA 2022
ALPHA_INV = 1 / ALPHA
LAMBDA = 4 * pi**3 + pi**2 + pi
D_LAMBDA = fabs(LAMBDA - ALPHA_INV)
recs = []
R = range(-10, 11)                   # 系数界 |a_i| <= 10
for a0 in R:
    for a1 in R:
        for a2 in R:
            for a3 in R:
                if (a0, a1, a2, a3) == (0, 0, 0, 0): continue
                v = a0 + a1*pi + a2*pi**2 + a3*pi**3
                recs.append((float(fabs(v - ALPHA_INV)), a0,a1,a2,a3))
recs.sort()
# 结果: 总候选 194480; Lambda(0,1,1,4) 排名 0 (全局最优);
# 次优 (-4,9,2,3) 偏差 3.626e-3 (差 12 倍); <1e-3: 1 个; <1e-2: 3 个; <0.1: 68 个
# 显著性判定见 q2_lambda_null_model.py (显式零模型, p ~ 0.13)
\end{lstlisting}

\begin{lstlisting}[language=Python,caption={Lambda 零模型与 $p$ 值（q2 / 定理 4 事后选择偏差检验）}]
# q2_lambda_null_model.py (节选; 完整见 src/)
# H0: 靶值 T 与格点取值集合无关, T ~ U[Vmin, Vmax]
# 统计量 D_min(T) = min_i |V_i - T|;  观察值 d_obs = |Lambda - alpha^{-1}|
V = np.sort(family_values)                      # 194480 个格点值
idx = np.searchsorted(V, T)                     # 精确最近邻, O(log N)
Dmin = np.minimum(np.abs(V[idx] - T), np.abs(V[idx-1] - T))
p_global = np.mean(Dmin <= d_obs)               # = 0.127 (4e5 次 MC)
# 解析: 1-(1-2d/R)^N = 0.123;  Bonferroni: N*2d/R = 0.132
# 限制 T in [100,200]: p = 0.173;  条件于 0.1 邻域 68 候选: p = 0.340
# 顺序统计量: Dmin 中位数 1.53e-3, 5% 分位 1.16e-4  => d_obs 仅居 12.7 百分位
# 置换检验退化性: 子集最佳拟合恒 >= d_obs, P(<=d_obs)=0.585 ~ 1-e^-1 (抽中 Lambda)
\end{lstlisting}

\begin{lstlisting}[language=Python,caption={Arb 认证区间（q3 / 定理 2、4、6 的严格判定）}]
# q3_arb_certified.py (节选; 完整见 src/)
from flint import arb, ctx          # Arb 球算术: 每步定向舍入扩展半径
ctx.prec = 128
alpha = arb("7.2973525643e-3", "1.1e-12")     # CODATA 2022 中心值 +/- u
mu0   = arb("1.25663706127e-6", "2.0e-16")
G     = arb("6.67430e-11", "1.5e-15")
a2mu0 = alpha * alpha * mu0
assert not a2mu0.overlaps(G)                  # 认证: 两区间不相交
assert a2mu0.lower() > G.upper()              # 差的下界 > 0
Lam  = 4*arb.pi()**3 + arb.pi()**2 + arb.pi()
assert not (Lam - 1/alpha).contains(arb(0))   # 认证: Lambda != alpha^{-1}
a2L2 = alpha*alpha * Lam*Lam
assert not a2L2.contains(arb(1))              # 认证: alpha^2 Lambda^2 != 1
# 输出: alpha^2 mu0 - G = [1.73125e-13, 1.76126e-13]（外取）; 间距 = 115.42 u(G); dev = 116.4 sigma
#       (注: str(arb) 的显示精度低于端点, 报告一律取 lower()/upper() 实值)
\end{lstlisting}

\begin{lstlisting}[language=Python,caption={mpmath 区间交叉核对（R4 / 与 Arb 结论一致性检查）}]
# R4_interval_arithmetic.py (节选; 完整见 src/)
from mpmath import mpi, mpf, pi
# CODATA 2022 中心值 +/- 标准不确定度 -> 区间
alpha = mpi("7.2973525643e-3", "1.1e-12")   # u = 1.1e-12 (u_r ~ 1.5e-10)
mu0   = mpi("1.25663706127e-6", "2.0e-16")  # 展示末位 2.0e-16；反除 u_r = 1.59e-10，\text{解析传递值} = u_r(alpha) = 1.51e-10
G     = mpi("6.67430e-11", "1.5e-15")       # u = 1.5e-15 (u_r ~ 2.2e-5)
a2mu0 = alpha**2 * mu0
# 结果: a2mu0 in [6.69176255e-11, 6.69176256e-11]
#       G     in [6.67415e-11, 6.67445e-11]
#       disjoint = True; 间距下界 1.73e-13
#       => alpha^2*mu0 > G 为区间可证的严格事实
# 注意: mpi 为误差传播包络, 非逐舍入认证 (区别于 Arb/MPFR)
\end{lstlisting}

\noindent\textbf{可复现性说明}：环境为 Python 3.13.7 + mpmath 1.3.0 + numpy 2.4.3 + python-flint 0.9.0（Arb 2.x），Lean 4.33.1；随机实验固定种子。全部工件已公开发布（见 §3 可得性段的仓库地址与许可）。

\section{2019 SI 下 $\mu_0$ 与 $\varepsilon_0$ 的导出关系与不确定度传递}

本附录补齐 §2 只给出结论的计量链条。

\textbf{推导}。精细结构常数的 SI 定义为
\[
  \alpha \;=\; \frac{e^2}{4\pi \varepsilon_0 \hbar c}.
\]
代入 $\varepsilon_0 \equiv 1/(\mu_0 c^2)$（这是 $c^2 = 1/(\mu_0\varepsilon_0)$ 的同一内容）得
\[
  \alpha = \frac{e^2 \mu_0 c^2}{4\pi \hbar c} = \frac{e^2 \mu_0 c}{4\pi\hbar}
  \;=\; \frac{e^2 \mu_0 c}{2 h},
\]
最后一步用 $\hbar = h/2\pi$。解出
\begin{equation}
  \boxed{\;\mu_0 \;=\; \frac{2\alpha h}{e^2 c}\;},\qquad
  \varepsilon_0 \;=\; \frac{1}{\mu_0 c^2} \;=\; \frac{e^2}{2\alpha h c}.
\end{equation}

\textbf{为什么 2019 年后 $\mu_0$ 不再是精确值}。2019 年 SI 重新定义把 $c$、$h$、$e$（以及 $\pi$、$4\pi$ 这类纯数）固定为精确值，而 $\alpha$ \textbf{不是}定义值——它由电子磁矩 $g$ 因子的量子比测定实验给出，是测量量。上式右端除 $\alpha$ 外全为精确值，故 $\mu_0$ 的全部不确定度来自 $\alpha$。

\textbf{不确定度传递}。设 $x_i$ 相互独立，$Y=\prod x_i^{p_i}$，则 $u_r(Y)=\sqrt{\sum (p_i u_r(x_i))^2}$。对 $\mu_0 = 2\alpha h e^{-2} c^{-1}$，指数为 $(p_\alpha,p_h,p_e,p_c)=(1,1,-2,-1)$，而 $h,e,c$ 精确（$u_r=0$），于是
\begin{equation}
  u_r(\mu_0) \;=\; u_r(\alpha), \qquad
  u_r(\varepsilon_0) \;=\; u_r(\mu_0) \;=\; u_r(\alpha)
\end{equation}
（$\varepsilon_0$ 一步用到 $c$ 精确，故 $u_r(\varepsilon_0)=u_r(\mu_0)$）。代入 CODATA 2022 的 $u(\alpha)=1.1\times10^{-12}$：
\begin{align}
  u_r(\alpha) &= 1.1\times10^{-12}/7.2973525643\times10^{-3} = 1.51\times10^{-10},\\
  u(\mu_0) &= u_r(\alpha)\cdot\mu_0 = 1.51\times10^{-10}\times1.25663706127\times10^{-6} = 1.9\times10^{-16},
\end{align}
与 CODATA 给出的 $\mu_0 = 1.256\,637\,061\,27(20)\times10^{-6}$ 同量级（$1.894\to2.0\times10^{-16}$ 是向上舍入到两位展示，相差 $5.6\%$——这 $5.6\%$ 正是表~\ref{tab:constants} 注\textsuperscript{a} 里“两口径反除不相等”的来源）。同理 $u(\varepsilon_0) = 1.51\times10^{-10}\times 8.854\,187\,8188\times10^{-12} = 1.34\times10^{-21}$。NIST 在线库把 $\varepsilon_0$ 展示为 $8.854\,187\,8188(14)\times10^{-12}$~F/m（未舍入导出值为 $8.854\,187\,818\,79\ldots\times10^{-12}$）。这里有一个可以\textbf{反过来用}的线索：上式给的 $u(\varepsilon_0)=1.34\times10^{-21}$ 只会舍入成 $(13)$，而 NIST 展示的是 $(14)$，说明未舍入的公共 $u_r$ 比本文按 $(11)$ 反除的 $1.51\times10^{-10}$ \textbf{还要大}（要大到能把 $u(\varepsilon_0)$ 推过 $1.35\times10^{-21}$、把 $u(\mu_0)$ 推过 $1.95\times10^{-16}$，即 $u_r\ge1.552\times10^{-10}$）。这正是表~\ref{tab:constants} 注\textsuperscript{a} 的交集论证所用的证据：本文 $\alpha$ 端球偏窄 2.9\%--4.5\%，其影响已由三档压力测试（\path|src/q5b_alpha_ball_nist_caliber.py|）覆盖——判定不变。

\textbf{一个易错点（本文自查中实际发生）}。符号 $7.2973525643(11)\times10^{-3}$ 的括号数字作用在\textbf{末位}，故 $u(\alpha)=1.1\times10^{-12}$ 而非 $1.1\times10^{-13}$；把绝对不确定度取小一个数量级会让 $\alpha^2\mu_0$ 的区间宽度虚减 10 倍。本文 §4.8 的区间已按正确取值重算；由于 $G$ 端不确定度比 $\alpha$ 端大五个数量级，不相交性与 $116\sigma$ 的结论不受此修正影响。

\textbf{对本文主结论的影响}。定理~1 的恒等式 $c^2\varepsilon_0\mu_0=1$ 是\textbf{定义层面}的关系，在 2019 SI 下依然成立（$\varepsilon_0$ 就是由 $\mu_0$ 与 $c$ 定义的导出量）；变化的只是"$\mu_0$ 是精确值"这一表述。由于 $u_r(\alpha^2\mu_0)= 3u_r(\alpha)\approx 4.5\times10^{-10} \ll u_r(G)=2.2\times10^{-5}$，$\alpha^2\mu_0$ 一侧的精度远高于 $G$ 一侧，$0.2616\%$ 偏差与 $116\sigma$ 的判定不因 SI 重新定义而改变。

\section{量纲格形式化的定理清单}

以下定理全部在 \texttt{G\_dimension\_lattice.lean} 中给出 Lean 4 证明（不依赖 Mathlib，仅用 \texttt{Int} 与 \texttt{omega}），经 \texttt{\#print axioms} 核验无 \texttt{sorryAx}。

\begin{center}
\footnotesize
\begin{tabular}{@{}llp{7.6cm}@{}}
\toprule
\textbf{Lean 名} & \textbf{对应正文} & \textbf{内容} \\
\midrule
\texttt{dimMono\_eq\_zero\_iff} & §4.7(四) & 单项式 $G^a\mu_0^c c^d\hbar^f e^g m^h$ 无量纲 $\iff$ 四条线性丢番图方程 \\
\texttt{dimMono\_param} & §4.7(四)1 & 无量纲指数格秩 2，生成元 $\Pi_1$ 与 $\mu_0ce^2/\hbar=4\pi\alpha$ \\
\texttt{no\_dimensionless\_G} & §4.7(四)2 & 无质量标度（$h=0$）时 $a=0$：$G$ 不进入任何无量纲单项式 \\
\texttt{Pi1\_dimensionless} & §4.7(二) & $\Pi_1 = Gm^2/(\hbar c)$ 无量纲（机器验证） \\
\texttt{alpha\_dimensionless} & §4.7(四)3 & $\alpha=\mu_0ce^2/(4\pi\hbar)$ 无量纲（机器验证） \\
\texttt{dimG\_ne\_dim\_mu0} & 定理~3 & $[G]\neq[\mu_0]=[\alpha^2\mu_0]$：SI 量纲矛盾可判定 \\
\texttt{cert\_gap\_pos} & §4.8 & Arb 证书间距 $>0$（整数标度 $10^{-22}$） \\
\texttt{alpha2mu0\_gt\_G\_of\_cert} & §4.8 & 证书回放：$x\ge$ 下界 $\wedge$ $y\le$ 上界 $\Rightarrow x>y$ \\
\bottomrule
\end{tabular}
\end{center}

\noindent 形式化的\textbf{边界}须同时声明：上述结果是\textbf{单项式}（幂积）层面的不变量分类。把它推广到含超越函数的任意表达式，需要"量纲不变性对函数复合保持"这类额外假设，本文不作此推广，也不声称已机械化"任何形式的可能关系都不存在"。

\end{document}
